\documentclass[11pt,a4paper]{report}

% --- Packages de base ---
\usepackage[utf8]{inputenc}
\usepackage[T1]{fontenc}
\usepackage[french]{babel} % Typo française
\usepackage{graphicx}
\usepackage{geometry}
\geometry{a4paper, top=3cm, bottom=3cm, left=2.5cm, right=2.5cm}
\usepackage{hyperref}
\usepackage{fancyhdr}
\usepackage{amsmath, amssymb, amsthm}
\usepackage{booktabs}
\usepackage{microtype}
\usepackage{float}
\usepackage{algorithm}
\usepackage{algpseudocode}
\floatname{algorithm}{Procédure}

\theoremstyle{plain}
\newtheorem{theorem}{Théorème}[chapter]
\newtheorem{lemma}[theorem]{Lemme}
\newtheorem{corollary}[theorem]{Corollaire}
\newtheorem{proposition}[theorem]{Proposition}
\theoremstyle{definition}
\newtheorem{definition}[theorem]{Définition}
\renewcommand{\proofname}{Preuve}

% --- Packages pour le design ---
\usepackage{eso-pic}       
\usepackage{tikz}          
\usepackage{titlesec}      

% --- Configuration des liens ---
\hypersetup{
    colorlinks=true,
    linkcolor=black,
    filecolor=magenta,      
    urlcolor=blue,
    pdftitle={Travail de Bachelor - Elie Bussod},
}

% --- Configuration des en-têtes et pieds de page (fancyhdr) ---
\pagestyle{fancy}
\fancyhf{} 
\fancyhead[L]{\leftmark} 
\fancyhead[R]{\thepage}  
\renewcommand{\headrulewidth}{0.4pt} 
\renewcommand{\footrulewidth}{0pt}   

% --- Configuration des gros numéros de chapitre (titlesec) ---
\titleformat{\chapter}[display]
  {\normalfont\bfseries\filright} 
  {\fontsize{80}{90}\selectfont\color{gray}\thechapter} 
  {1ex} 
  {\Huge\color{black}} 
\titlespacing*{\chapter}{0pt}{10pt}{40pt}

\begin{document}

% =============================================================

% PAGE DE GARDE
\begin{titlepage}
    % écusson en fond 
    \AddToShipoutPictureBG*{
        \begin{tikzpicture}[remember picture, overlay]
            % position
            \node[anchor=south east, opacity=0.2, rotate=25] 
            at ([xshift=3.1cm, yshift=4cm]current page.south east) {
                
\includegraphics[width=1\textwidth]{images/ecusson_unige.png}
            };
        \end{tikzpicture}
    }
    
    % logos
    \noindent
    \begin{minipage}[c]{0.45\textwidth}
        \begin{flushleft}
            
\includegraphics[height=1.75cm]{images/logo_unige.png}
        \end{flushleft}
    \end{minipage}
    \hfill
    \begin{minipage}[c]{0.45\textwidth}
        \begin{flushright}
            
\includegraphics[height=2.9cm]{images/logo_unige_sciences.png}
        \end{flushright}
    \end{minipage}
    
    \begin{center}
        \vspace*{1.5cm}
        
        {\Large \textsc{Université de Genève}} \\[0.5cm]
        {\large \textsc{Faculté des Sciences}} \\[0.2cm]
        {\large Département d'Informatique} \\[2cm]
        
        {\large Travail de Bachelor en Sciences Informatiques} \\[1cm]
        {\large Année académique 2025-2026} \\[2cm]
        
        % Titre
        \rule{\linewidth}{0.5mm} \\[0.5cm]
        {\huge \bfseries Programmation dynamique pour formules SAT \\[0.4cm] }
        \rule{\linewidth}{0.5mm} \\[2.5cm]
        
    \end{center}
    
    % Auteur et Superviseurs
    \begin{flushleft} \large
        \textbf{Auteur :}\\
        Elie \textsc{Bussod} \\[1cm]
        \textbf{Superviseurs :} \\
        Prof. Arnaud \textsc{Casteigts} \\
        Prof. Pierre \textsc{Leone}
    \end{flushleft}
    
    \vfill
    \begin{center}
        {\large Date de rendu : Juin 2026}
    \end{center}
\end{titlepage}

% =============================================================


% Résumé / Abstract 
\chapter*{Résumé}
\addcontentsline{toc}{chapter}{Résumé}

Nous implémentons en C le comptage de modèles (\#SAT) et la
maximisation de clauses (MaxSAT) par programmation dynamique sur
branch decomposition, dont la complexité est paramétrée par la
ps-width (algorithme de Saether, Telle et Vatshelle). L'implémentation
est étendue de manière à produire, au cours de la même exécution, un
DAG d-DNNF structuré (construction de Bova, Capelli, Mengel et
Slivovsky), sur lequel six requêtes polytimes sont évaluées.
L'évaluation empirique montre qu'un solveur de production reste plus
rapide en résolution unique, mais que le DAG compilé s'amortit dès
qu'une même formule est interrogée de nombreuses fois, ce qui place la
contribution du côté de la compilation de connaissances plutôt que de
la résolution.

\clearpage

% Table des matières
\tableofcontents
\clearpage

% Liste des figures
\listoffigures
\clearpage


% Introduction
\chapter{Introduction}
\label{chap:intro}

\section{Contexte et motivation}

La satisfaisabilité propositionnelle (SAT) occupe une place centrale en informatique
théorique et appliquée depuis Cook \cite{cook1971}, qui en a établi la
NP-complétude, et Karp \cite{karp1972}, qui a tiré de cette propriété une vaste
famille de réductions polynomiales et fait de SAT le problème de référence de la
théorie de la complexité. Malgré cette intractabilité théorique dans le cas général, les
solveurs modernes parviennent à traiter des instances industrielles comportant des
millions de variables et de clauses, grâce notamment à l'algorithme d'apprentissage
de clauses guidé par les conflits (CDCL) \cite{marques1999,een2003}, à la propagation
unitaire (BCP) et à des heuristiques de branchement sophistiquées. Cette efficacité
pratique a conduit à une adoption massive de SAT dans des domaines aussi variés que
la vérification formelle de circuits électroniques, la planification automatique, la
synthèse de programmes et la cryptanalyse.

Ce succès a naturellement suscité un intérêt croissant pour des variantes plus
expressives : le \emph{comptage de modèles} (\#SAT) et la \emph{maximisation de
clauses satisfaites} (MaxSAT). Le problème \#SAT consiste à déterminer le nombre
exact d'affectations satisfaisant une formule donnée en forme normale
conjonctive (CNF). Il généralise la
décision de satisfaisabilité en quantifiant la cardinalité de l'ensemble des
modèles, et intervient naturellement dans le raisonnement probabiliste,
l'inférence dans les réseaux bayésiens et la vérification probabiliste
\cite{darwiche2002}. Le problème MaxSAT, quant à lui, cherche une affectation
qui satisfait le plus grand nombre possible de clauses, même lorsque la formule
est globalement insatisfaisable. Il modélise des problèmes d'optimisation
combinatoire à contraintes souples (ordonnancement, configuration, bioinformatique)
et constitue un cadre unifiant pour de nombreux problèmes NP-difficiles.

Ces deux variantes sont strictement plus difficiles que SAT : \#SAT est
\#P-complet \cite{valiant1979} et MaxSAT est NP-difficile. Cette complexité
accrue justifie le développement d'approches algorithmiques spécialisées.

Une stratégie complémentaire à la résolution directe consiste à transformer
hors-ligne la formule en une représentation structurée, typiquement un graphe
orienté acyclique de type d-DNNF, sur laquelle un large éventail de requêtes
(comptage, satisfaisabilité, équivalence) deviennent polynomiales. Cette
\emph{compilation de connaissances} \cite{darwiche2002} amortit le coût
hors-ligne sur de multiples interrogations ultérieures, et est devenue le
standard pour l'inférence probabiliste exacte et la vérification quantitative.
Elle constitue le second pilier de ce mémoire, aux côtés de la résolution
paramétrée par la structure de la formule.

\section{Approche structurelle et décomposition arborescente}

Plusieurs familles d'algorithmes ont émergé pour résoudre \#SAT et MaxSAT :
solveurs généraux fondés sur DPLL (Davis-Putnam-Logemann-Loveland) et CDCL,
outils SMT (\emph{satisfiability modulo theories}) comme \texttt{Z3}
\cite{moura2008}, solveurs de comptage avec cache de composantes comme
\texttt{sharpSAT} \cite{thurley2006}, ou encore la compilation de connaissances
déjà évoquée. Ce travail combine deux perspectives complémentaires : la
\emph{complexité paramétrée structurelle}, qui exploite la structure du graphe
d'incidence (graphe biparti reliant chaque variable aux clauses où elle
apparaît) pour résoudre \#SAT et MaxSAT en temps polynomial dès que le
paramètre est borné ; et la \emph{compilation de connaissances}, qui transforme
cette même structure en un objet polynomialement interrogeable.

L'idée fondatrice est que la difficulté d'une instance ne dépend pas uniquement
de sa taille $n + m$ (variables et clauses), mais de paramètres structurels
liés à la façon dont variables et clauses interagissent. Si ces paramètres
sont petits, des algorithmes par \emph{programmation dynamique sur décomposition
arborescente} résolvent \#SAT et MaxSAT en temps polynomial en $n + m$,
l'explosion exponentielle étant confinée à ce seul paramètre.

Parmi ces paramètres, la \emph{treewidth} du graphe d'incidence de la formule
est le plus étudié \cite{samer2010,slivovsky2020}. Des résultats plus récents,
fondés sur la \emph{mim-width} \cite{vatshelle2012} et la \emph{ps-width}
\cite{saether2015}, étendent la tractabilité à des classes de formules
strictement plus larges. L'algorithme de Saether, Telle et Vatshelle
\cite{saether2015} résout simultanément \#SAT et MaxSAT en temps
$O(k^3 \cdot m \cdot (m+n))$, où $k$ est la ps-width de la décomposition.
Il constitue le premier pilier de ce mémoire de bachelor. Son extension par Bova, Capelli,
Mengel et Slivovsky \cite{bova2015} en forme le second : ces auteurs
l'instrumentent pour produire en parallèle un DAG d-DNNF, faisant ainsi le pont
avec la compilation de connaissances.

\section{Objectifs de ce travail}

Ce mémoire de bachelor poursuit deux objectifs complémentaires.

Le premier est l'\textbf{implémentation complète} de l'algorithme de Saether,
Telle et Vatshelle \cite{saether2015} en langage C. Cette implémentation comprend
le parseur de formules DIMACS CNF, trois stratégies de construction de l'arbre
de décomposition (dont l'heuristique GreedyOrder décrite dans le papier), les
trois procédures de l'algorithme (calcul bottom-up des PS-sets, propagation
top-down, programmation dynamique), ainsi que les structures de données
sous-jacentes (bitsets 64-bit, trie binaire, tables DP).

Le second objectif est l'\textbf{extension à la compilation de connaissances}
fondée sur le résultat de Bova, Capelli, Mengel et Slivovsky
\cite{bova2015}. Leur algorithme s'obtient comme une instrumentation de la
programmation dynamique précédente et produit, à coût marginal, un DAG
d-DNNF de taille polynomiale en la ps-width et la taille de la formule. Sur le DAG produit,
l'ensemble des requêtes que la carte de compilation de Darwiche et Marquis
\cite{darwiche2002} rend polynomiales sur d-DNNF est implémenté.

Une \textbf{évaluation empirique} accompagne ces deux volets, sur des
instances synthétiques conçues pour contrôler la ps-width, des formules
aléatoires $k$-SAT et des circuits de multiplication encodés par la
transformation de Tseytin. Ces tests mesurent l'impact de la ps-width sur
les performances et délimitent les régimes dans lesquels la décomposition
structurelle apporte un avantage décisif.

% =============================================================
% ===========================================================================
\chapter{État de l'art}
\label{chap:etat_art}
% ===========================================================================

Ce chapitre positionne notre travail dans la littérature. Nous présentons
successivement les problèmes SAT, \#SAT et MaxSAT, les principales familles
d'algorithmes de résolution, les paramètres de décomposition structurelle
qui gouvernent leur complexité, l'algorithme de Saether, Telle et Vatshelle
sur lequel repose notre implémentation, et la compilation de connaissances
CNF $\to$ d-DNNF par Bova, Capelli, Mengel et Slivovsky qui en constitue
l'extension naturelle.

% ---------------------------------------------------------------------------
\section{SAT, \#SAT et MaxSAT}
% ---------------------------------------------------------------------------

\subsection{Le problème SAT}

Le problème de satisfaisabilité propositionnelle (SAT) demande, étant donnée
une formule $F$ en forme normale conjonctive (CNF), s'il existe une
affectation des variables qui la satisfait. Rappelons qu'une formule CNF est
une conjonction de clauses, chaque clause étant une disjonction de littéraux
(variable ou sa négation). Une affectation satisfaisante est appelée
\emph{modèle} de $F$.

SAT fut le premier problème formellement démontré NP-complet par Cook
\cite{cook1971}, puis Karp en étendit l'applicabilité en montrant que de
nombreux problèmes combinatoires classiques s'y réduisent en temps
polynomial \cite{karp1972}. Malgré cette intractabilité théorique, les
solveurs industriels modernes traitent couramment des instances de plusieurs
millions de variables et clauses, grâce à l'algorithme
CDCL \cite{moskewicz2001,een2003} décrit à la section \ref{sec:cdcl}.

\subsection{\#SAT : comptage de modèles}

Le problème \#SAT (\emph{model counting}) demande non plus si un modèle
existe, mais \emph{combien} il en existe. Ce problème est \#P-complet
\cite{valiant1979}, classe de complexité qui capture le comptage associé
aux décisions NP. Le comptage reste \#P-complet même sur des fragments
syntaxiques pour lesquels la décision est facile (2-CNF, Horn).

\#SAT intervient dans l'inférence probabiliste sur des réseaux bayésiens
\cite{darwiche2002}, le model checking quantitatif et l'analyse de fiabilité
de systèmes. Sa résolution exacte demande d'explorer l'intégralité de
l'espace des affectations.

\subsection{MaxSAT}

Le problème MaxSAT cherche une affectation qui maximise le nombre de clauses
satisfaites, même si la formule est globalement insatisfaisable. Il est
NP-difficile et modélise des problèmes d'optimisation à contraintes souples
en ordonnancement, configuration ou bioinformatique. Sa variante \emph{Weighted
MaxSAT}, où chaque clause porte un poids et l'on maximise la somme des poids
satisfaits, est celle que résout l'algorithme étudié dans ce mémoire.

% ---------------------------------------------------------------------------
\section{Algorithmes de résolution}
\label{sec:algos}
% ---------------------------------------------------------------------------

\subsection{DPLL et CDCL}
\label{sec:cdcl}

L'architecture CDCL (\emph{Conflict-Driven Clause Learning}) étend
l'algorithme historique DPLL (Davis-Putnam-Logemann-Loveland) par
l'\emph{apprentissage de clauses} : à chaque conflit, le solveur extrait
une nouvelle clause qui résume la cause de l'échec et qu'il ajoute à la
formule. Popularisée par zChaff \cite{moskewicz2001} puis MiniSAT
\cite{een2003}, CDCL est aujourd'hui le standard industriel et traite
couramment des instances de plusieurs millions de variables.

\subsection{Extensions pour \#SAT et MaxSAT}

Pour le comptage de modèles, les solveurs étendent DPLL par le \emph{cache
de composantes} : lorsque l'affectation partielle déconnecte le graphe
d'incidence, les comptes des composantes disjointes sont calculés
indépendamment puis multipliés. Ce mécanisme, popularisé par
\texttt{sharpSAT} \cite{thurley2006}, est repris par \texttt{D4}
\cite{lagniez2017} qui compile en plus la formule vers d-DNNF.
Pour la maximisation, les solveurs MaxSAT modernes reposent sur des
approches \emph{core-guided}, qui résolvent itérativement des noyaux
insatisfaisables.

\subsection{Compilation de connaissances}

Une approche orthogonale à la résolution en ligne, la \emph{compilation de
connaissances}, transforme hors-ligne la formule en une représentation
polynomialement interrogeable et fait l'objet de la
section \ref{sec:kc}.

Toutes ces approches ont en commun de ne pas exploiter directement la
structure topologique de la formule. La section suivante présente une
famille d'algorithmes qui en fait le paramètre central.

% ---------------------------------------------------------------------------
\section{Décompositions arborescentes et paramètres de largeur}
\label{sec:decomp}
% ---------------------------------------------------------------------------

Pour les formules issues du monde réel, la difficulté pratique dépend
moins de la taille brute que de la structure de connectivité du
\emph{graphe d'incidence} de la formule : il s'agit du graphe biparti
$I(F)$ dont les sommets sont d'un côté les variables $\mathsf{var}(F)$,
de l'autre les clauses $\mathsf{cla}(F)$, avec une arête entre une
variable $x$ et une clause $C$ chaque fois que $x$ apparaît dans $C$.
Plusieurs paramètres ont été proposés pour quantifier cette structure et
fournir des algorithmes de programmation dynamique tractables quand le
paramètre est borné.

Le plus étudié est la \emph{treewidth}, introduite par Robertson et
Seymour dans leur théorie des mineurs de graphes : on représente le
graphe par une famille de sous-ensembles de sommets (\emph{sacs})
connectés en arbre, et la largeur de la décomposition est la taille
du plus grand sac diminuée de un. Samer et Szeider \cite{samer2010}
puis Slivovsky et Szeider \cite{slivovsky2020} ont obtenu des algorithmes
de DP pour \#SAT paramétrés par la treewidth du graphe d'incidence.
Toutefois, la treewidth n'est bornée que sur des graphes globalement
clairsemés, ce qui en limite la portée sur les classes denses.

Pour dépasser cette limitation, des paramètres plus généraux ont été
introduits, dont la \emph{mim-width} (\emph{maximum induced
matching width}) de Vatshelle \cite{vatshelle2012}. Elle se définit
via une décomposition arborescente du graphe : à chaque coupe, on
mesure la taille du plus grand \emph{couplage induit} traversant
cette coupe (un couplage induit étant un ensemble d'arêtes deux à deux
non adjacentes, sans autre arête du graphe entre leurs extrémités) ;
la mim-width est le minimum sur toutes les décompositions du maximum
de cette quantité. Elle reste bornée sur des classes à treewidth non
bornée (graphes d'intervalles ou de permutation, par exemple). En
contrepartie, calculer la mim-width est intractable : Bergougnoux,
Bonnet et Duron \cite{bergougnoux2025} ont montré que le problème est
paraNP-complet, c'est-à-dire NP-difficile même pour des graphes de
mim-width bornée par une constante. Cela impose le recours à des
heuristiques dans toute mise en oeuvre pratique.

% ---------------------------------------------------------------------------
\section{ps-width et l'algorithme de Saether-Telle-Vatshelle}
\label{sec:pswidth}
% ---------------------------------------------------------------------------

Saether, Telle et Vatshelle \cite{saether2015} introduisent la
\emph{ps-width} (\emph{projection-satisfiability width}), un paramètre
défini directement sur les formules CNF à partir d'une \emph{branch
decomposition} de la formule. Il domine la mim-width du graphe
d'incidence : pour des clauses de taille bornée, une mim-width bornée
entraîne une ps-width polynomiale en le nombre de clauses, ce qui étend
strictement les classes tractables au-delà de la treewidth. Les auteurs
établissent que \#SAT et MaxSAT se résolvent par programmation dynamique
sur la décomposition en temps $O(k^3 \cdot m \cdot (m+n))$, ramené à
$O(k^2 \cdot m \cdot (m+n))$ dans le cas linéaire.

Comme minimiser la ps-width est intractable, ils proposent l'heuristique
gloutonne \emph{GreedyOrder} et valident l'algorithme sur trois familles
synthétiques d'instances à ps-width contrôlée : \texttt{sharpSAT} et les
solveurs MaxSAT modernes restent plus rapides sur la plupart des entrées,
mais sur les familles construites pour favoriser le DP leur prototype les
bat nettement \cite{saether2015}. La définition formelle du paramètre,
la description des trois procédures de l'algorithme et le fonctionnement
de l'heuristique font l'objet du chapitre \ref{chap:theorie} ; les
familles d'instances servent de banc de test au
chapitre \ref{chap:resultats}.

% ---------------------------------------------------------------------------
\section{Compilation de connaissances CNF \texorpdfstring{$\to$}{vers} d-DNNF}
\label{sec:kc}
% ---------------------------------------------------------------------------

La \emph{compilation de connaissances} (\emph{knowledge compilation}, KC)
transforme une formule CNF hors-ligne en une représentation structurée,
typiquement un graphe orienté acyclique (DAG), sur laquelle certaines
requêtes (comptage, conditionnement, projection partielle) deviennent
polynomiales. Le coût
de compilation est amorti dès qu'une même base est interrogée plusieurs
fois ; la compilation est aujourd'hui un outil standard de l'inférence
probabiliste exacte. Cette section présente la hiérarchie des langages
cibles, les requêtes tractables sur d-DNNF, puis le résultat de Bova,
Capelli, Mengel et Slivovsky (BCMS 2015) sur lequel repose l'extension
KC de notre implémentation.

\subsection{Langages cibles}

Darwiche et Marquis \cite{darwiche2002} ont dressé la \emph{carte de
compilation} qui classe les représentations selon deux axes : la
\emph{succinctitude} (taille relative) et la \emph{tractabilité} des
requêtes supportées en temps polynomial. La hiérarchie pertinente pour
notre travail va de la NNF (\emph{Negation Normal Form} : DAG dont les
feuilles sont des littéraux et les noeuds internes des ET et OU) à la
DNNF (\emph{Decomposable NNF} : enfants de tout ET à portées variables
disjointes), puis au \textbf{d-DNNF} (\emph{Deterministic DNNF} : enfants
de tout OU à ensembles de modèles disjoints). Cette dernière propriété
rend le comptage \#SAT polynomial par un parcours qui remplace OU par
somme et ET par produit, et fait du d-DNNF le langage cible de notre
implémentation. Le SDD (\emph{Sentential Decision Diagram}
\cite{darwiche2011}) est un d-DNNF dont la structure est dictée par un
\emph{vtree}, un arbre binaire qui fixe l'ordre des variables aux
feuilles. L'OBDD (\emph{Ordered Binary Decision Diagram}
\cite{bryant1986}) est le précurseur historique. Ces deux variantes
sont plus restrictives que d-DNNF et ne sont pas utilisées dans ce
travail.

Au sein de la classe d-DNNF, deux sous-langages distincts se sont
développés en parallèle. Les compilateurs DPLL-based, dont \texttt{c2d}
puis \texttt{D4} \cite{lagniez2017}, produisent des
\emph{decision-DNNF} via une recherche CDCL avec cache de composantes :
les noeuds OU correspondent à des branchements sur une variable et la
taille du DAG résultant n'est pas bornée a priori en termes de paramètres
structurels de la formule. Les \emph{d-DNNF structurés}, à l'inverse,
imposent un schéma global (typiquement une branch decomposition fixée
en amont) qui structure les noeuds AND et permet d'établir des bornes
théoriques sur la taille de la représentation. La construction de
Bova, Capelli, Mengel et Slivovsky (section \ref{subsec:bcms}) appartient
à cette seconde famille et fournit une borne théorique sur la taille du
DAG paramétrée par la ps-width.

\subsection{Requêtes tractables sur d-DNNF}
\label{subsec:queries-ddnnf}

Darwiche et Marquis \cite{darwiche2002} établissent que la classe d-DNNF
rend polynomiales six requêtes standard : \textbf{CO}
(\emph{consistency}), \textbf{VA} (\emph{validity}), \textbf{CE}
(\emph{clausal entailment}), \textbf{IM} (\emph{implicant check}),
\textbf{CT} (\emph{model counting}, le \#SAT) et \textbf{ME}
(\emph{model enumeration}, énumération avec délai polynomial entre
modèles consécutifs). Le DAG produit par la construction de la
section \ref{subsec:bcms} satisfait les hypothèses requises et ces
six requêtes y sont donc directement disponibles. Leur algorithmique
précise est décrite au chapitre \ref{chap:theorie}.

\subsection{Compilation explicite pilotée par la ps-width : BCMS 2015}
\label{subsec:bcms}

Bova, Capelli, Mengel et Slivovsky \cite{bova2015} prolongent
l'algorithme de Saether, Telle et Vatshelle en montrant que toute formule
CNF $F$ dotée d'une branch decomposition de ps-width $k$ peut être
compilée en un d-DNNF structuré de taille
\[
  |D| \leq 7 \cdot k^3 \cdot (n + m),
\]
en temps polynomial en $n + m$. À la différence des compilateurs
DPLL-based (\texttt{c2d}, \texttt{D4}), la construction n'effectue
aucune recherche : la décomposition fournit le squelette du DAG, le DP
fournit son contenu, et la borne $7 \cdot k^3 \cdot (n+m)$ donne une
garantie de taille théorique paramétrée par la ps-width, absente des
compilateurs top-down. La construction détaillée du DAG est décrite au
chapitre \ref{chap:theorie}.


\chapter{Cadre théorique}
\label{chap:theorie}

Les sections \ref{sec:theo-prelim} et \ref{sec:theo-pswidth} fixent les
notations et définissent la branch decomposition puis la ps-width. La
section \ref{sec:theo-dp} expose la programmation dynamique de Saether,
Telle et Vatshelle \cite{saether2015} qui résout \#SAT et MaxSAT en
temps polynomial pour ps-width bornée, et la section \ref{sec:theo-greedy}
l'heuristique gloutonne GreedyOrder qui lui fournit une décomposition.
Bova, Capelli, Mengel et Slivovsky \cite{bova2015}
(section \ref{sec:theo-bcms}) instrumentent ce DP pour produire un
d-DNNF structuré, sur lequel la section \ref{sec:theo-queries} décrit
les six requêtes polytimes.

% -----------------------------------------------------------------------------
\section{Notations et préliminaires}
\label{sec:theo-prelim}
% -----------------------------------------------------------------------------

On considère une formule $F$ en CNF, vue comme un multi-ensemble de
clauses, chaque clause étant un ensemble de littéraux (une variable
ou sa négation). On note $\mathsf{var}(F)$ l'ensemble des variables
apparaissant dans $F$ et $\mathsf{cla}(F)$ son ensemble de clauses.
Une affectation d'un ensemble de variables $X$ est une fonction
$\tau : X \to \{0,1\}$, étendue aux littéraux par
$\tau(\neg x) = 1 - \tau(x)$. Elle satisfait une clause $C$ s'il
existe un littéral $\ell \in C$ tel que $\tau(\ell) = 1$, et satisfait
la formule $F$ si elle satisfait toutes les clauses de
$\mathsf{cla}(F)$.

Pour $\mathcal{C} \subseteq \mathsf{cla}(F)$ et
$X \subseteq \mathsf{var}(F)$, la \emph{sous-formule induite}
$F_{\mathcal{C},X}$ s'obtient en retirant de $F$ toute clause hors de
$\mathcal{C}$, puis tout littéral dont la variable n'appartient pas
à $X$. Cette construction décrira les sous-formules de part et
d'autre d'une coupe de la décomposition (section \ref{sec:theo-pswidth}).

Le graphe d'incidence $I(F)$ a été défini au
chapitre \ref{chap:etat_art}, section \ref{sec:decomp}. Ses sommets
sont les variables et clauses de $F$, avec une arête entre $x$ et $C$
si et seulement si $x$ apparaît dans $C$. On suit l'hypothèse de
Saether, Telle et Vatshelle \cite{saether2015} en supposant $I(F)$
connexe : sinon le problème se résout indépendamment sur chaque
composante connexe.

% -----------------------------------------------------------------------------
\section{Branch decomposition et ps-width}
\label{sec:theo-pswidth}
% -----------------------------------------------------------------------------

\subsection{Branch decomposition d'une formule}
\label{ssec:theo-branch-decomp}

Une \emph{branch decomposition} d'une formule $F$ est un couple
$(T, \delta)$ où $T$ est un arbre binaire enraciné et $\delta$ une
bijection entre les feuilles de $T$ et l'ensemble
$\mathsf{var}(F) \cup \mathsf{cla}(F)$ \cite{saether2015}. Chaque
variable et chaque clause de $F$ apparaît ainsi exactement une fois
comme étiquette d'une feuille. Pour un noeud non-feuille $v$ de $T$,
on note $\delta(v)$ l'ensemble des étiquettes des feuilles du
sous-arbre enraciné en $v$ :
\[
  \delta(v) = \{\, \delta(l) : l \text{ feuille du sous-arbre enraciné en } v\,\}.
\]
Le noeud $v$ partage $\mathsf{var}(F) \cup \mathsf{cla}(F)$ en deux,
$\delta(v)$ et son complément. Cette partition est appelée la
\emph{coupe} associée à $v$.

À chaque noeud $v$ sont associées deux sous-formules. En écrivant
$\delta(v) = \mathcal{C} \cup X$ avec
$\mathcal{C} \subseteq \mathsf{cla}(F)$ et $X \subseteq \mathsf{var}(F)$,
on pose
\[
  F_v = F_{\overline{\mathcal{C}}, X}
  \qquad \text{et} \qquad
  F_{\overline{v}} = F_{\mathcal{C}, \overline{X}},
\]
où $\overline{\mathcal{C}} = \mathsf{cla}(F) \setminus \mathcal{C}$ et
$\overline{X} = \mathsf{var}(F) \setminus X$, et où la notation
$F_{\mathcal{C}, X}$ a été fixée en section \ref{sec:theo-prelim}.
L'asymétrie est volontaire : la sous-formule $F_v$ rassemble les
clauses extérieures à la coupe restreintes aux variables locales, et
décrit précisément ce qu'une affectation des seules variables de $X$
peut satisfaire parmi les clauses du complément.

\begin{figure}[htbp]
\centering
\begin{tikzpicture}[
  level distance=10mm,
  level 1/.style={sibling distance=58mm},
  level 2/.style={sibling distance=22mm},
  level 3/.style={sibling distance=11mm},
  every node/.style={font=\small},
  internal/.style={circle, fill=black, inner sep=1pt},
  vnode/.style={circle, draw=black, thick, fill=white,
                inner sep=1pt, minimum size=5mm, font=\itshape\small},
  leaf/.style={draw, rectangle, rounded corners=2pt, fill=gray!10,
               inner sep=3pt, minimum height=4mm, font=\small},
]
\node[internal] {}
  child {node[vnode] {$v$}
    child {node[internal] {}
      child {node[leaf] {$x_1$}}
      child {node[leaf] {$x_2$}}
    }
    child {node[internal] {}
      child {node[leaf] {$c_1$}}
      child {node[leaf] {$c_3$}}
    }
  }
  child {node[internal] {}
    child {node[internal] {}
      child {node[leaf] {$x_3$}}
      child {node[leaf] {$c_4$}}
    }
    child {node[internal] {}
      child {node[internal] {}
        child {node[leaf] {$x_5$}}
        child {node[leaf] {$c_2$}}
      }
      child {node[leaf] {$x_4$}}
    }
  };
\end{tikzpicture}
\caption{Une branch decomposition $(T, \delta)$ d'une formule $F$
sur cinq variables et quatre clauses, reprise de la Figure 2 de
\cite{saether2015}. La formule a pour variables
$\mathsf{var}(F) = \{x_1, x_2, x_3, x_4, x_5\}$ et pour clauses
$c_1 = \{x_1, x_2\}$, $c_2 = \{x_1, \neg x_2, x_3\}$,
$c_3 = \{x_2, \neg x_4, x_5\}$, $c_4 = \{x_2, x_3, x_5\}$. Le
noeud $v$ marqué définit la coupe
$\delta(v) = \{x_1, x_2, c_1, c_3\}$, soit
$\mathcal{C} = \{c_1, c_3\}$ et $X = \{x_1, x_2\}$. Les deux
sous-formules croisées sont
$F_v = \{\{x_1, \neg x_2\}, \{x_2\}\}$ (versions de $c_2$ et $c_4$
restreintes aux variables de $X$) et
$F_{\overline{v}} = \{\emptyset, \{\neg x_4, x_5\}\}$ (versions de
$c_1$ et $c_3$ restreintes aux variables de $\overline{X}$).}
\label{fig:branch-decomp-example}
\end{figure}

Quand tous les noeuds internes de $T$ se trouvent sur un même chemin
(autrement dit, chaque noeud interne a au moins une feuille pour
enfant), on parle de \emph{décomposition linéaire}. Ce cas particulier
admet une analyse de complexité plus fine, traitée en
section \ref{ssec:theo-th23}.

\subsection{Coupes et PS-sets}
\label{ssec:theo-cuts-psset}

Soient $\mathcal{C} \subseteq \mathsf{cla}(F)$ et
$X \subseteq \mathsf{var}(F)$. L'ensemble $S = \mathcal{C} \cup X$ est
une \emph{coupe} de $F$. Chaque noeud non-feuille $v$ d'une branch
decomposition en définit une via $\delta(v)$. Une coupe sépare $F$ en
quatre sous-formules induites $F_{\mathcal{C},X}$,
$F_{\overline{\mathcal{C}},X}$, $F_{\mathcal{C},\overline{X}}$ et
$F_{\overline{\mathcal{C}},\overline{X}}$, chacune obtenue en
restreignant clauses et variables aux sous-ensembles qui composent $S$
ou son complément.

Pour une affectation $\tau$ d'un sous-ensemble de variables, on note
$\mathsf{sat}(F, \tau)$ l'ensemble inclusion-maximal des clauses de $F$
satisfaites par $\tau$. Un sous-ensemble
$\mathcal{C} \subseteq \mathsf{cla}(F)$ est dit
\emph{projection-satisfiable} dans $F$ s'il existe une affectation
$\tau$ portant sur l'ensemble complet $\mathsf{var}(F)$ telle que
$\mathsf{sat}(F, \tau) = \mathcal{C}$. On dit alors que $\mathcal{C}$
est une \emph{projection} de $F$ par $\tau$.

La famille des sous-ensembles projection-satisfiables de $F$ se note
\[
  \mathsf{PS}(F) = \{\, \mathsf{sat}(F, \tau) : \tau \text{ est une affectation de } \mathsf{var}(F)\,\}.
\]
Sa cardinalité $|\mathsf{PS}(F)|$ est la \emph{ps-value} de $F$. C'est
à partir de cette quantité, calculée sur les sous-formules $F_v$ et
$F_{\overline{v}}$ associées aux noeuds d'une branch decomposition, que
la section \ref{ssec:theo-pswidth-def} définira la ps-width.

\subsection{ps-value et ps-width}
\label{ssec:theo-pswidth-def}

Pour un noeud $v$ d'une branch decomposition $(T, \delta)$, la
\emph{ps-value} de la coupe associée se définit par
\[
  \mathsf{ps}(\delta(v)) = \max\bigl(|\mathsf{PS}(F_v)|,\
  |\mathsf{PS}(F_{\overline{v}})|\bigr).
\]
La \emph{ps-width} de la décomposition est la plus grande de ces
ps-values sur tous les noeuds de $T$ :
\[
  \mathsf{psw}(T, \delta) = \max\,\bigl\{\,\mathsf{ps}(\delta(v)) :
  v \in T\,\bigr\}.
\]
La \emph{ps-width} de la formule $F$ est le minimum atteint sur
toutes ses branch decompositions :
\[
  \mathsf{psw}(F) = \min\,\bigl\{\,\mathsf{psw}(T, \delta) : (T, \delta)
  \text{ branch decomposition de } F\,\bigr\}.
\]

Pour une clause $c$ et un ensemble de variables $Y$, on note $c|_Y$
la clause obtenue en ne gardant que les littéraux dont la variable est
dans $Y$ : $c|_Y = \{\, \ell \in c : \mathsf{var}(\ell) \in Y\,\}$.
Détaillons le calcul de $|\mathsf{PS}(F_v)|$ sur l'exemple de la
Figure \ref{fig:branch-decomp-example}. La sous-formule $F_v$ contient
les versions restreintes de $c_2$ et $c_4$ aux variables de $X = \{x_1, x_2\}$ :
\[
  D_1 = c_2|_X = \{x_1, \neg x_2\},
  \qquad
  D_2 = c_4|_X = \{x_2\}.
\]
Les quatre affectations possibles sur $\mathsf{var}(F_v) = \{x_1, x_2\}$
donnent :

\begin{center}
\begin{tabular}{cc|cc|c}
\hline
$x_1$ & $x_2$ & $D_1$ & $D_2$ & $\mathsf{sat}(F_v, \tau)$ \\
\hline
0 & 0 & satisfaite & non satisfaite & $\{D_1\}$ \\
0 & 1 & non satisfaite & satisfaite & $\{D_2\}$ \\
1 & 0 & satisfaite & non satisfaite & $\{D_1\}$ \\
1 & 1 & satisfaite & satisfaite & $\{D_1, D_2\}$ \\
\hline
\end{tabular}
\end{center}

Les quatre affectations produisent trois ensembles satisfaits distincts,
d'où $\mathsf{PS}(F_v) = \{\{D_1\}, \\ \{D_2\}, \{D_1, D_2\}\}$ et
$|\mathsf{PS}(F_v)| = 3$. Du côté de
$F_{\overline{v}} = \{\emptyset, \{\neg x_4, x_5\}\}$, la clause vide
n'est jamais satisfaite, et les affectations sur $\{x_4, x_5\}$ qui
satisfont $\{\neg x_4, x_5\}$ donnent
$\mathsf{sat} = \{c_3|_{\overline{X}}\}$, les autres donnent
$\mathsf{sat} = \emptyset$. D'où
$\mathsf{PS}(F_{\overline{v}}) = \{\emptyset, \{c_3|_{\overline{X}}\}\}$
et $|\mathsf{PS}(F_{\overline{v}})| = 2$. La ps-value de la coupe vaut
donc $\mathsf{ps}(\delta(v)) = \max(3, 2) = 3$.

La programmation dynamique de la section \ref{sec:theo-dp} manipule
plus naturellement les clauses originales de $F$ que leurs versions
restreintes. Le papier introduit pour cela une variante de notation :
$\mathsf{PS}'(F_v)$ désigne la famille des sous-ensembles
$\mathcal{D} \subseteq \mathsf{cla}(F) \setminus \delta(v)$ dont la
restriction aux variables de $F_v$ appartient à $\mathsf{PS}(F_v)$, et
$\mathsf{sat}'(F_v, \tau)$ est définie de manière analogue à partir de
$\mathsf{sat}(F_v, \tau)$. Sur notre exemple,
$\mathsf{PS}'(F_v) = \{\{c_2\}, \{c_4\}, \{c_2, c_4\}\}$ : chaque
élément de $\mathsf{PS}(F_v)$ est remplacé par le sous-ensemble
correspondant de clauses originales de $F$. La correspondance entre
$\mathsf{PS}$ et $\mathsf{PS}'$ est bijective :
$|\mathsf{PS}'(F_v)| = |\mathsf{PS}(F_v)|$ et la ps-value est
inchangée.

Saether, Telle et Vatshelle établissent par ailleurs (Lemme 1 de
\cite{saether2015}) que pour des clauses de taille bornée,
$|\mathsf{PS}(F)|$ reste polynomial en le nombre de clauses dès que la
taille maximale d'un couplage induit du graphe d'incidence $I(F)$ est
bornée : la ps-width est donc un paramètre tractable sur les classes
de formules dont la mim-width (définie en section \ref{sec:decomp})
est bornée.

% -----------------------------------------------------------------------------
\section{Programmation dynamique pour \#SAT et MaxSAT}
\label{sec:theo-dp}
% -----------------------------------------------------------------------------

\subsection{Stratégie d'ensemble}
\label{ssec:theo-dp-strategy}

Le calcul de \#SAT et MaxSAT exploite la structure récursive de la
branch decomposition $(T, \delta)$. À chaque noeud $v$ de $T$, la coupe
$\delta(v)$ délimite les sous-formules $F_v$ et $F_{\overline{v}}$, et
toute affectation $\tau$ de $\mathsf{var}(F)$ se décompose en sa
restriction $\tau|_v$ aux variables de $F_v$ et sa restriction
$\tau|_{\overline{v}}$ aux variables de $F_{\overline{v}}$. Comme une
clause peut être satisfaite par des littéraux situés des deux côtés de
la coupe, la programmation dynamique de Saether, Telle et
Vatshelle \cite{saether2015} retient à chaque coupe l'ensemble des
clauses déjà satisfaites par chaque moitié, ce qui permet de combiner
les contributions des deux côtés sans énumérer toutes les
affectations. Les PS-sets introduits en
section \ref{ssec:theo-pswidth-def} fournissent l'espace d'états
adapté : la ps-width $k$ borne leur cardinal à chaque coupe, et les
boucles internes du DP restent polynomiales en $k$.

L'algorithme s'organise en trois procédures successives. La
\emph{Procédure 1} calcule, par parcours bottom-up de $T$, les
ensembles $\mathsf{PS}'(F_v)$. La \emph{Procédure 2} calcule, par
parcours top-down, les ensembles duaux
$\mathsf{PS}'(F_{\overline{v}})$. Ces deux passes sont des
pré-traitements obligatoires : la \emph{Procédure 3}, qui effectue le
DP proprement dit en remontant l'arbre, a besoin en chaque noeud des
PS-sets des deux côtés de la coupe pour indexer ses tables. La valeur
cherchée se lit ensuite à la racine, où la table se réduit à une
unique cellule. Les sections \ref{ssec:theo-proc1}
à \ref{ssec:theo-th23} en formalisent successivement chaque étape.

\subsection{Procédure 1 : calcul bottom-up des PS'(F\textsubscript{v})}
\label{ssec:theo-proc1}

La Procédure 1 calcule $\mathsf{PS}'(F_v)$ pour tous les noeuds de la
branch decomposition par parcours bottom-up de $T$. À chaque noeud, le
résultat dépend uniquement des PS-sets calculés aux deux enfants, et
la récursion atteint les feuilles, où le calcul est explicite.

Aux feuilles, $\mathsf{var}(F_\ell)$ contient au plus une variable et
$\mathsf{PS}'(F_\ell)$ se construit directement. Si $\ell$ est étiquetée
par une variable $x$, les deux affectations $\tau(x) = 1$ et
$\tau(x) = 0$ produisent
\[
  \mathsf{sat}'(F_\ell, \tau{=}1) = \{c \in \mathsf{cla}(F) : x \in c\},
  \qquad
  \mathsf{sat}'(F_\ell, \tau{=}0) = \{c \in \mathsf{cla}(F) : \neg x \in c\},
\]
qui sont les deux éléments candidats de $\mathsf{PS}'(F_\ell)$. Le
cardinal de $\mathsf{PS}'(F_\ell)$ vaut $1$ ou $2$ selon que ces deux
ensembles coïncident ou non. Si $\ell$ est étiquetée par une clause
$c$, alors $\mathsf{var}(F_\ell) = \emptyset$, l'unique affectation
$\varepsilon$ (sur l'ensemble vide de variables) ne satisfait aucune
clause restreinte à $\emptyset$, et $\mathsf{PS}'(F_\ell) =
\{\emptyset\}$.

Pour un noeud interne $v$ d'enfants $c_1$ et $c_2$, le PS-set se
reconstruit récursivement par
\[
  \mathsf{PS}'(F_v) = \bigl\{\,(C_1 \cup C_2) \cap \mathsf{cla}(F_v) :
  C_1 \in \mathsf{PS}'(F_{c_1}),\ C_2 \in \mathsf{PS}'(F_{c_2})\,\bigr\}.
\]
La formule reflète la décomposition $\mathsf{var}(F_v) =
\mathsf{var}(F_{c_1}) \sqcup \mathsf{var}(F_{c_2})$ : une affectation
$\tau$ de $\mathsf{var}(F_v)$ se sépare en $\tau_1 =
\tau|_{\mathsf{var}(F_{c_1})}$ et $\tau_2 = \tau|_{\mathsf{var}(F_{c_2})}$,
et une clause $c \in \mathsf{cla}(F_v)$ est satisfaite par $\tau$ si et
seulement si $\tau_1$ ou $\tau_2$ y satisfait un littéral. L'intersection
avec $\mathsf{cla}(F_v) = \mathsf{cla}(F) \setminus \delta(v)$ écarte les
clauses étiquetant les feuilles du sous-arbre, qui sortent du domaine
de $F_v$.

\begin{algorithm}[H]
\caption{Génération de $\mathsf{PS}'(F_v)$}
\label{algo:proc1}
\begin{algorithmic}[1]
\Require $\mathsf{PS}'(F_{c_1})$, $\mathsf{PS}'(F_{c_2})$ pour les enfants $c_1, c_2$ de $v$
\Ensure $\mathsf{PS}'(F_v)$
\State $L \leftarrow \emptyset$
\ForAll{$(C_1, C_2) \in \mathsf{PS}'(F_{c_1}) \times \mathsf{PS}'(F_{c_2})$}
  \State ajouter $(C_1 \cup C_2) \cap \mathsf{cla}(F_v)$ à $L$
\EndFor
\State \Return $L$
\end{algorithmic}
\end{algorithm}

\begin{theorem}[Théorème 1 de \cite{saether2015}]
\label{thm:proc1-complexity}
Soit $F$ une formule CNF avec $n = |\mathsf{var}(F)|$ et
$m = |\mathsf{cla}(F)|$, et soit $(T, \delta)$ une branch decomposition
de $F$ de ps-width $k$. L'ensemble des $\mathsf{PS}'(F_v)$ et des
$\mathsf{PS}'(F_{\overline{v}})$ pour tous les noeuds $v$ de $T$ se calcule
en temps $O(k^2\, m\, (m + n))$.
\end{theorem}

\subsection{Procédure 2 : calcul top-down des PS'(F\textsubscript{$\overline{v}$})}
\label{ssec:theo-proc2}

La Procédure 2 calcule $\mathsf{PS}'(F_{\overline{v}})$ pour tous les
noeuds $v$ de $T$ par parcours top-down. Elle remplit le rôle symétrique
de la Procédure 1 sur l'autre côté de chaque coupe : la récursion
descend de la racine vers les feuilles, et chaque calcul à un noeud
non-racine combine deux PS-sets déjà disponibles, celui du frère $s$
produit par la Procédure 1 et celui du parent $p$ obtenu plus haut dans
la descente.

À la racine $r$, $\mathsf{var}(F_{\overline{r}}) = \emptyset$ et l'unique
affectation $\varepsilon$ ne satisfait aucune clause restreinte à
$\emptyset$. On a donc $\mathsf{PS}(F_{\overline{r}}) = \{\emptyset\}$,
d'où le cas de base
\[
  \mathsf{PS}'(F_{\overline{r}}) = \{\emptyset\}.
\]

Pour un noeud non-racine $v$ de frère $s$ et de parent $p$, le PS-set
s'obtient par
\[
  \mathsf{PS}'(F_{\overline{v}}) = \bigl\{\,(C_s \cup C_p) \cap \mathsf{cla}(F_{\overline{v}}) :
  C_s \in \mathsf{PS}'(F_s),\ C_p \in \mathsf{PS}'(F_{\overline{p}})\,\bigr\}.
\]
La formule reflète la décomposition $\mathsf{var}(F_{\overline{v}}) =
\mathsf{var}(F_s) \sqcup \mathsf{var}(F_{\overline{p}})$ : les variables
hors de $\delta(v)$ proviennent soit du sous-arbre du frère
($\mathsf{var}(F_s) \subseteq \delta(s)$), soit de l'extérieur du
sous-arbre du parent ($\mathsf{var}(F_{\overline{p}}) \subseteq
\overline{\delta(p)}$). Une affectation de $\mathsf{var}(F_{\overline{v}})$
se sépare donc en deux moitiés indépendantes, et l'intersection avec
$\mathsf{cla}(F_{\overline{v}}) = \mathsf{cla}(F) \cap \delta(v)$ écarte
les clauses étiquetant des feuilles extérieures au sous-arbre de $v$,
qui sortent du domaine de $F_{\overline{v}}$.

\begin{algorithm}[H]
\caption{Génération de $\mathsf{PS}'(F_{\overline{v}})$}
\label{algo:proc2}
\begin{algorithmic}[1]
\Require $\mathsf{PS}'(F_s)$, $\mathsf{PS}'(F_{\overline{p}})$ pour le frère $s$ et le parent $p$ de $v$
\Ensure $\mathsf{PS}'(F_{\overline{v}})$
\State $L \leftarrow \emptyset$
\ForAll{$(C_s, C_p) \in \mathsf{PS}'(F_s) \times \mathsf{PS}'(F_{\overline{p}})$}
  \State ajouter $(C_s \cup C_p) \cap \mathsf{cla}(F_{\overline{v}})$ à $L$
\EndFor
\State \Return $L$
\end{algorithmic}
\end{algorithm}

La Procédure 2 a la même structure que la Procédure 1 et s'analyse de
la même manière, d'où l'énoncé conjoint du
Théorème \ref{thm:proc1-complexity}.

\subsection{Procédure 3 : table DP $\mathsf{Tab}_v(C, C')$}
\label{ssec:theo-proc3}

La Procédure 3 calcule la table de programmation dynamique en chaque
noeud par parcours bottom-up de $T$. C'est l'étape qui exploite les
PS-sets produits par les Procédures 1 et 2 pour obtenir, à la racine,
la solution de MaxSAT et de \#SAT.

Soit $v$ un noeud de $T$ et
$(C, C') \in \mathsf{PS}'(F_v) \times \mathsf{PS}'(F_{\overline{v}})$. Une
affectation $\tau$ de $\mathsf{var}(F)$ \emph{rencontre l'expectation}
$(C, C')$ si elle vérifie simultanément
\begin{description}
  \item[(A)] $\mathsf{sat}'(F_v, \tau|_{\mathsf{var}(F_v)}) = C$,
  \item[(B)] $\mathsf{sat}'(F_{\overline{v}}, \tau|_{\mathsf{var}(F_{\overline{v}})}) = C'$.
\end{description}
La condition (A) fixe la trace que les variables intérieures à
$\delta(v)$ laissent sur les clauses extérieures, et la condition (B)
fixe la trace symétrique des variables extérieures sur les clauses
intérieures. La table de programmation dynamique attachée à $v$ est
l'application
\[
  \mathsf{Tab}_v : \mathsf{PS}'(F_v) \times \mathsf{PS}'(F_{\overline{v}}) \to \mathbb{N},
\]
où $\mathsf{Tab}_v(C, C')$ est le nombre maximal de clauses de $\delta(v)$
satisfaites par une affectation rencontrant l'expectation $(C, C')$.
Comme les variables hors de $\delta(v)$ y satisfont exactement les
clauses de $C'$, cette définition équivaut à la Constraint (1)
de \cite{saether2015} :
\begin{equation}
  \mathsf{Tab}_v(C, C') = \max_{\tau} \bigl|\,(\mathsf{sat}'(F, \tau) \cap \delta(v)) \cup C'\,\bigr|,
  \label{eq:tab-constraint}
\end{equation}
où $\tau$ parcourt les affectations de $\mathsf{var}(F) \cap \delta(v)$
telles que $\mathsf{sat}'(F_v, \tau) = C$.

Pour un noeud interne $v$ d'enfants $c_1$ et $c_2$, la valeur de
$\mathsf{Tab}_v(C_v, C'_v)$ se reconstruit à partir des tables filles
$\mathsf{Tab}_{c_1}$ et $\mathsf{Tab}_{c_2}$. Chaque table fille étant
indexée par un couple, une mise à jour fait intervenir six indices au
total. Une astuce due à Bui-Xuan, Telle et Vatshelle réduit la boucle
de six à trois indices : on itère sur les triplets
$(C_{c_1}, C_{c_2}, C'_v) \in \mathsf{PS}'(F_{c_1}) \times \mathsf{PS}'(F_{c_2}) \times \mathsf{PS}'(F_{\overline{v}})$,
et l'on dérive canoniquement
\[
  C'_{c_1} = (C_{c_2} \cup C'_v) \cap \delta(c_1),
  \qquad
  C'_{c_2} = (C_{c_1} \cup C'_v) \cap \delta(c_2),
  \qquad
  C_v = (C_{c_1} \cup C_{c_2}) \setminus \delta(v).
\]
Ces trois identités expriment les relations entre les couples
d'expectation qu'une même affectation rencontre aux coupes de $c_1$,
$c_2$ et $v$. La cellule $\mathsf{Tab}_v(C_v, C'_v)$ se met alors à jour
par le maximum des sommes
$\mathsf{Tab}_{c_1}(C_{c_1}, C'_{c_1}) + \mathsf{Tab}_{c_2}(C_{c_2}, C'_{c_2})$
ainsi obtenues.

\begin{algorithm}[H]
\caption{Calcul de $\mathsf{Tab}_v$ pour un noeud interne $v$ d'enfants $c_1, c_2$ (MaxSAT)}
\label{algo:proc3}
\begin{algorithmic}[1]
\Require $\mathsf{Tab}_{c_1}$, $\mathsf{Tab}_{c_2}$
\Ensure $\mathsf{Tab}_v$
\State $\mathsf{Tab}_v(C, C') \leftarrow -1$ pour tout $(C, C') \in \mathsf{PS}'(F_v) \times \mathsf{PS}'(F_{\overline{v}})$
\ForAll{$(C_{c_1}, C_{c_2}, C'_v) \in \mathsf{PS}'(F_{c_1}) \times \mathsf{PS}'(F_{c_2}) \times \mathsf{PS}'(F_{\overline{v}})$}
  \State $C'_{c_1} \leftarrow (C_{c_2} \cup C'_v) \cap \delta(c_1)$
  \State $C'_{c_2} \leftarrow (C_{c_1} \cup C'_v) \cap \delta(c_2)$
  \State $C_v \leftarrow (C_{c_1} \cup C_{c_2}) \setminus \delta(v)$
  \State $t \leftarrow \mathsf{Tab}_{c_1}(C_{c_1}, C'_{c_1}) + \mathsf{Tab}_{c_2}(C_{c_2}, C'_{c_2})$
  \If{$\mathsf{Tab}_v(C_v, C'_v) < t$}
    \State $\mathsf{Tab}_v(C_v, C'_v) \leftarrow t$
  \EndIf
\EndFor
\State \Return $\mathsf{Tab}_v$
\end{algorithmic}
\end{algorithm}

Pour \#SAT, la cellule $\mathsf{Tab}_v(C, C')$ compte plutôt le nombre
d'affectations $\tau$ de $\mathsf{var}(F) \cap \delta(v)$ telles que
$\mathsf{sat}'(F_v, \tau) = C$ et que toute clause de $\delta(v)$ soit
satisfaite par $\tau$ ou figure dans $C'$. La Procédure 3 s'adapte de
la manière suivante : la table s'initialise à $0$ partout, les feuilles
reçoivent les valeurs $0$, $1$ ou $2$ données par la définition, et les
lignes 6 et 7 du pseudocode sont remplacées par
\[
  \mathsf{Tab}_v(C_v, C'_v) \leftarrow \mathsf{Tab}_v(C_v, C'_v)
  + \mathsf{Tab}_{c_1}(C_{c_1}, C'_{c_1}) \cdot \mathsf{Tab}_{c_2}(C_{c_2}, C'_{c_2}).
\]
La somme cumulée remplace le maximum, et le produit des deux tables
filles dénombre les combinaisons compatibles à chaque itération.

\begin{lemma}[Lemme 2 de \cite{saether2015}]
\label{lem:proc3-complexity}
Soit $F$ une formule CNF de $m$ clauses et $v$ un noeud interne d'une
branch decomposition $(T, \delta)$ de ps-width $k$. La Procédure 3
calcule $\mathsf{Tab}_v$ vérifiant l'équation \eqref{eq:tab-constraint}
en temps $O(k^3\, m)$.
\end{lemma}

\subsection{Théorèmes 2 et 3}
\label{ssec:theo-th23}

Une fois les trois procédures exécutées sur $T$, le résultat du DP se
lit en une seule cellule à la racine. Comme noté en
section \ref{ssec:theo-dp-strategy}, à la racine $r$ on a
$\mathsf{PS}'(F_r) = \mathsf{PS}'(F_{\overline{r}}) = \{\emptyset\}$, et
la table $\mathsf{Tab}_r$ se réduit à l'unique cellule
$\mathsf{Tab}_r(\emptyset, \emptyset)$.

\begin{theorem}[Théorème 2 de \cite{saether2015}]
\label{thm:dp-total}
Soit $F$ une formule CNF de $n$ variables et $m$ clauses, et
$(T, \delta)$ une branch decomposition de $F$ de ps-width $k$. Le DP
calcule MaxSAT et \#SAT en temps $O(k^3\, m\, (m + n))$, la valeur
cherchée se lisant à la racine en $\mathsf{Tab}_r(\emptyset, \emptyset)$.
\end{theorem}

La validité vient de l'équation \eqref{eq:tab-constraint} : à la racine,
$\mathsf{cla}(F) \cap \delta(r) = \mathsf{cla}(F)$, donc
$\mathsf{Tab}_r(\emptyset, \emptyset)$ vaut le maximum du nombre de
clauses satisfiables par une affectation de $\mathsf{var}(F)$ (MaxSAT).
La variante additive de la Procédure 3 produit, au même endroit, le
nombre d'affectations qui satisfont toutes les clauses (\#SAT). La
complexité s'obtient en sommant le coût du pré-traitement
(Théorème \ref{thm:proc1-complexity}, $O(k^2\, m\, (m + n))$) et celui
de la Procédure 3 sur les $O(m + n)$ noeuds internes
(Lemme \ref{lem:proc3-complexity}, $O(k^3\, m)$ par noeud).

Le facteur $k^3$ provient du triplet sur lequel boucle la Procédure 3 à
chaque noeud interne. Sur une décomposition \emph{linéaire} (chaque noeud
interne a au moins un enfant feuille), le PS-set de cet enfant est de
cardinal au plus $2$, soit constant. Le triplet se réduit donc à
$O(k^2)$ termes, et la borne globale tombe d'un facteur $k$.

\begin{theorem}[Théorème 3 de \cite{saether2015}]
\label{thm:dp-linear}
Sous les hypothèses du Théorème \ref{thm:dp-total}, si $(T, \delta)$ est
une décomposition linéaire, alors MaxSAT et \#SAT se résolvent en temps
$O(k^2\, m\, (m + n))$.
\end{theorem}

Cette amélioration ne demande aucune modification de l'algorithme : la
boucle de la Procédure 3 itère sur les tailles effectives des PS-sets
des enfants, et l'un d'eux étant constant sur une décomposition
linéaire, le compte d'itérations par noeud passe à $O(k^2)$. Les
décompositions produites par l'heuristique GreedyOrder, présentée en
section \ref{sec:theo-greedy}, sont précisément de ce type.

% -----------------------------------------------------------------------------
\section{Heuristique GreedyOrder}
\label{sec:theo-greedy}
% -----------------------------------------------------------------------------

L'algorithme du DP exposé en section \ref{sec:theo-dp} prend en entrée
une branch decomposition $(T, \delta)$. Pour minimiser le facteur $k$
dans la complexité, il faudrait disposer de la ps-width minimale parmi
toutes les branch decompositions de $F$. Calculer cette quantité est
probablement intractable \cite{saether2015}, mais aucune démonstration
formelle n'en a été donnée. En pratique, le DP s'exécute donc sur une
décomposition produite par une heuristique, sans garantie d'optimalité.

GreedyOrder est l'heuristique élémentaire proposée par le papier pour
cet usage. Elle prend en entrée le graphe d'incidence $I(F)$ et produit
un ordre linéaire $\sigma$ sur ses sommets (variables et clauses
confondues). À chaque étape, elle choisit comme prochain sommet celui
ayant le plus grand nombre de voisins \emph{déjà choisis}, en
départageant les ex aequo par le plus petit degré total.

\begin{algorithm}[H]
\caption{GreedyOrder}
\label{algo:greedyorder}
\begin{algorithmic}[1]
\Require Graphe biparti $G = (V, E)$
\Ensure Ordre linéaire $\sigma$ sur $V$
\State $L \leftarrow \emptyset$, $R \leftarrow V$, $i \leftarrow 1$
\ForAll{$v \in V$} \State $\mathit{Ldegree}(v) \leftarrow 0$ \EndFor
\While{$R \neq \emptyset$}
  \State choisir $v$ : parmi les sommets de $R$ de $\mathit{Ldegree}$ maximal, prendre un de plus petit degré total
  \State $\sigma(i) \leftarrow v$, $i \leftarrow i + 1$, transférer $v$ de $R$ vers $L$
  \ForAll{$w \in R$ tel que $vw \in E$}
    \State $\mathit{Ldegree}(w) \leftarrow \mathit{Ldegree}(w) + 1$
  \EndFor
\EndWhile
\State \Return $\sigma$
\end{algorithmic}
\end{algorithm}

L'ordre $\sigma$ se traduit en une branch decomposition linéaire
$(T, \delta)$ : $T$ est un arbre binaire dont les noeuds internes
forment un chemin enraciné à une extrémité, et $\delta$ envoie la
$i$-ème feuille rencontrée par un parcours en largeur partant de la
racine sur le sommet $\sigma(i)$. La structure résultante est un
\emph{caterpillar} : chaque noeud interne a une feuille pour enfant, à
l'exception du noeud le plus profond qui en a deux. C'est précisément
la forme d'arbre où la borne $O(k^2\, m\, (m+n))$ du
Théorème \ref{thm:dp-linear} s'applique.

GreedyOrder n'offre aucune garantie théorique sur la ps-width de la
décomposition produite : les auteurs la qualifient eux-mêmes de
« very simple heuristic » et la présentent comme une simple preuve de
concept. Son intérêt réside dans la simplicité de mise en oeuvre et dans
les performances empiriques observées au chapitre \ref{chap:resultats}.

% -----------------------------------------------------------------------------
\section{Compilation de F en d-DNNF structuré (BCMS)}
\label{sec:theo-bcms}
% -----------------------------------------------------------------------------

Bova, Capelli, Mengel et Slivovsky \cite{bova2015} (BCMS dans la suite)
montrent qu'à partir de la programmation dynamique de Saether, on peut
construire en parallèle un DAG d-DNNF qui calcule la formule $F$. Leur
papier introduit pour cela un objet appelé \emph{shape}, noté
$\mathbf{S} = (S, S')$. Comme cette notion coïncide avec les paires
$(C, C')$ déjà utilisées en section \ref{sec:theo-dp} (mêmes ensembles
sous-jacents, et même sémantique que la variante \#SAT du DP de
Saether, section \ref{ssec:theo-proc3}), nous conservons la notation
$(C, C')$ dans tout ce qui suit pour préserver l'unité de lecture du
chapitre. Les écarts ponctuels avec la nomenclature BCMS seront
signalés explicitement.

\subsection{Construction du DAG par récurrence}
\label{ssec:theo-bcms-construction}

BCMS associe à chaque noeud $v$ et chaque paire $(C, C') \in
\mathsf{PS}'(F_v) \times \mathsf{PS}'(F_{\overline{v}})$ un d-DNNF
$\varphi_v(C, C')$ construit récursivement à partir de ses enfants. La
récurrence repose sur les mêmes dérivations à trois indices que la
Procédure 3 (section \ref{ssec:theo-proc3}), vues comme une relation
de \emph{génération} entre une paire-parent et un couple de
paires-enfants : $(C_1, C_1')$ et $(C_2, C_2')$ \emph{génèrent}
$(C, C')$ si
\[
  C = (C_1 \cup C_2) \setminus \delta(v),
  \quad
  C_1' = (C' \cup C_2) \cap \delta(c_1),
  \quad
  C_2' = (C' \cup C_1) \cap \delta(c_2).
\]
On note $G(C, C')$ l'ensemble des couples qui génèrent $(C, C')$ en
$v$, et $N_v(C, C')$ l'ensemble des affectations
$\tau : \mathsf{var}(F_v) \to \{0, 1\}$ satisfaisant les conditions de
la variante \#SAT du DP (section \ref{ssec:theo-proc3}).

Aux feuilles de $T$, les ensembles $N_v(C, C')$ sont explicites et les
fragments $\varphi_v$ se réduisent à des littéraux ou des constantes.
Pour une \textbf{feuille variable} $\delta(v) = \{x\}$, soient
$\tau_x, \tau_{\neg x}$ les deux affectations de $\{x\}$ et
$C_\ell = \mathsf{sat}'(F_v, \tau_\ell)$ pour $\ell \in \{x, \neg x\}$.
Les paires admises sont $(C_x, \emptyset)$ et $(C_{\neg x}, \emptyset)$,
et l'on pose
\[
  \varphi_v(C_x, \emptyset) \equiv x,
  \qquad
  \varphi_v(C_{\neg x}, \emptyset) \equiv \neg x.
\]
Pour une \textbf{feuille clause} $\delta(v) = \{c\}$, l'unique
affectation $\varepsilon$ sur l'ensemble vide de variables ne satisfait
aucune clause. Les paires admises sont $(\emptyset, \emptyset)$ et
$(\emptyset, \{c\})$. Comme $\varepsilon \in N_v(\emptyset, \{c\})$
mais $\varepsilon \notin N_v(\emptyset, \emptyset)$, on pose
\[
  \varphi_v(\emptyset, \emptyset) \equiv \bot,
  \qquad
  \varphi_v(\emptyset, \{c\}) \equiv \top.
\]

Pour un noeud interne $v$ d'enfants $c_1$ et $c_2$, et une paire
$(C, C')$ en $v$, on combine les fragments des enfants via
\begin{equation}
  \varphi_v(C, C') \equiv
  \bigvee_{((C_1, C_1'), (C_2, C_2')) \in G(C, C')}
  \varphi_{c_1}(C_1, C_1') \wedge \varphi_{c_2}(C_2, C_2').
  \label{eq:bcms-construction}
\end{equation}
Un noeud AND conjugue les fragments des deux enfants pour chaque couple
générateur, et tous ces noeuds AND alimentent un même noeud OR qui forme
la racine de $\varphi_v(C, C')$. Les constantes $\bot$ et $\top$ sont
propagées en cours de construction et les sous-termes triviaux
supprimés.

Le d-DNNF $\varphi_v(C, C')$ ainsi obtenu a deux propriétés : ses
variables sont contenues dans $\mathsf{var}(F_v)$, et ses modèles sont
exactement les affectations de $N_v(C, C')$. C'est le contenu du
Lemme 4 de BCMS, énoncé en
section \ref{ssec:theo-bcms-properties}.

À la racine $r$ de $T$, l'unique paire est $(\emptyset, \emptyset)$, et
la compilation de $F$ s'obtient comme $\varphi_r(\emptyset, \emptyset)$
par évaluation de la récurrence \eqref{eq:bcms-construction} en
remontant l'arbre. La correction de cette compilation est formalisée
en section \ref{ssec:theo-bcms-properties}.

\subsection{Propriétés du DAG produit}
\label{ssec:theo-bcms-properties}

La construction décrite à la
section \ref{ssec:theo-bcms-construction} est correcte comme
compilateur de $F$ en d-DNNF structuré. Sa justification s'organise
autour de trois lemmes du papier BCMS, que l'on énonce sans
démonstration.

\begin{lemma}[Lemme 4 de \cite{bova2015}]
\label{lem:bcms-correctness}
Pour tout noeud $v$ de $T$ et toute paire $(C, C')$ admise en $v$, le
DAG $\varphi_v(C, C')$ est un d-DNNF dont les variables sont incluses
dans $\mathsf{var}(F_v)$. Une affectation
$\tau : \mathsf{var}(F_v) \to \{0,1\}$ satisfait $\varphi_v(C, C')$ si
et seulement si $\tau \in N_v(C, C')$.
\end{lemma}

Pour énoncer la structuration, on associe à chaque noeud $v$ un
\emph{vtree} de $\mathsf{var}(F_v)$, noté $\mathsf{vtree}(T, v)$. C'est
le sous-arbre $T_v$ dans lequel on supprime les feuilles étiquetées
par une clause, puis dans lequel on contracte les arêtes pendantes
pour restaurer la binarité (Figure 1 de \cite{bova2015}). Le résultat
est une branch decomposition de $\mathsf{var}(F_v)$ et joue le rôle
d'arbre de référence pour la structuration du d-DNNF.

\begin{lemma}[Lemme 5 de \cite{bova2015}]
\label{lem:bcms-structured}
Pour tout noeud $v$ et toute paire $(C, C')$ admise en $v$, le d-DNNF
$\varphi_v(C, C')$ respecte $\mathsf{vtree}(T, v)$. Le DAG construit
est donc un d-DNNF structuré.
\end{lemma}

\begin{lemma}[Lemme 6 de \cite{bova2015}]
\label{lem:bcms-root}
À la racine $r$ de $T$, la seule paire admise est
$(\emptyset, \emptyset)$, et $\varphi_r(\emptyset, \emptyset)$ est un
d-DNNF structuré dont les modèles sont exactement les affectations de
$\mathsf{var}(F)$ satisfaisant $F$.
\end{lemma}

Les trois lemmes ci-dessus assemblent un compilateur de $F$ en d-DNNF
structuré, sur lequel les requêtes polytimes étudiées en
section \ref{sec:theo-queries} sont directement applicables. La taille
du DAG produit fait l'objet de la section suivante.

\subsection{Borne de taille (Lemme 7 et Théorème 1)}
\label{ssec:theo-bcms-bound}

\begin{lemma}[Lemme 7 de \cite{bova2015}]
\label{lem:bcms-size}
Soit $F$ une formule CNF de $n$ variables et $m$ clauses, et
$(T, \delta)$ une branch decomposition de $F$ de ps-width $k$. Le
d-DNNF $\varphi_r(\emptyset, \emptyset)$ produit par la construction
de la section \ref{ssec:theo-bcms-construction} a une taille bornée
par
\[
  |\varphi_r(\emptyset, \emptyset)| \leq 7\, k^3\, (n + m).
\]
\end{lemma}

La borne se justifie par décompte. Chaque triplet $(C_1, C_2, C')$ de
la section \ref{ssec:theo-proc3} ajoute au plus cinq noeuds et arêtes
au DAG : un noeud AND, deux arêtes vers les fragments des enfants, une
arête vers le noeud OR de la paire-parent, et éventuellement un nouveau
noeud OR à allouer. Avec au plus $k^3$ triplets par noeud interne et
$n + m - 1$ noeuds internes dans $T$, on obtient $5 k^3 (n + m - 1)$
noeuds et arêtes au plus. Chaque feuille ajoute au plus deux fragments
mono-noeuds (littéraux ou constantes), soit $2 (n + m)$ au total. La
somme $5 k^3 (n + m - 1) + 2 (n + m)$ se majore par $7 k^3 (n + m)$
pour $k \geq 1$.

\begin{theorem}[Théorème 1 de \cite{bova2015}]
\label{thm:bcms-main}
Toute formule CNF $F$ de $n$ variables, $m$ clauses et de ps-width $k$
admet une compilation en d-DNNF structuré de taille
$O(k^3 (n + m))$. L'algorithme de compilation associé s'exécute en
temps $O(k^3\, m\, (n + m))$.
\end{theorem}

La borne en taille reprend le Lemme \ref{lem:bcms-size}. Le temps
incorpore le coût $O(m)$ par opération ensembliste, déjà présent dans
la complexité du DP de Saether (Théorème \ref{thm:dp-total}). La
validation empirique de cette borne sur les instances de test fait
l'objet du chapitre \ref{chap:resultats}.

Le d-DNNF produit appartient à la classe d-DNNF générale, plus large
que celle des decision-DNNF
(chapitre \ref{chap:etat_art}, section \ref{sec:kc}). Ce point n'a
pas d'incidence sur l'applicabilité des requêtes polytimes étudiées
au chapitre suivant.

% -----------------------------------------------------------------------------
\section{Requêtes polytimes sur d-DNNF}
\label{sec:theo-queries}
% -----------------------------------------------------------------------------

La compilation BCMS (section \ref{sec:theo-bcms}) produit un d-DNNF
qui sert de support à six requêtes calculables en temps polynomial,
déjà nommées au chapitre \ref{chap:etat_art},
section \ref{subsec:queries-ddnnf} : consistance (CO), validité (VA),
implication clausale (CE), test d'implicant (IM), comptage des
modèles (CT, c'est-à-dire \#SAT) et énumération des modèles (ME). Le
résultat fondateur est dû à Darwiche et Marquis \cite{darwiche2002},
qui établissent que chacune de ces requêtes se réduit à un parcours
du DAG dont la complexité est polynomiale en sa taille, dès que
celui-ci satisfait certaines propriétés structurelles
(sous-section \ref{ssec:theo-properties}).

\subsection{Trois propriétés clés}
\label{ssec:theo-properties}

Les trois propriétés clés du d-DNNF sont la \emph{décomposabilité}
(enfants d'un AND à portées variables disjointes), le
\emph{déterminisme} (enfants d'un OR à ensembles de modèles disjoints)
et la \emph{lissité} (enfants d'un OR à portées variables identiques).
Les deux premières définissent la classe d-DNNF, déjà introduite au
chapitre \ref{chap:etat_art}, section \ref{sec:kc}. La lissité, qui
n'y figure pas, simplifie plusieurs des requêtes en garantissant que
tous les sous-DAGs d'un OR parlent du même ensemble de variables.

Par le Lemme \ref{lem:bcms-correctness}, le DAG construit en
section \ref{sec:theo-bcms} est décomposable et déterministe. Il n'est
pas nécessairement lisse, mais peut être rendu tel à la demande par
l'opération de lissage décrite en sous-section \ref{ssec:theo-trans}.

\subsection{CT : model counting}
\label{ssec:theo-ct}

Le comptage des modèles d'un d-DNNF $\varphi$ se calcule par un
parcours post-ordre du DAG. La fonction $\mathsf{ct}(\varphi)$ se
définit récursivement par
\[
  \mathsf{ct}(\bot) = 0,
  \quad
  \mathsf{ct}(\top) = 1,
  \quad
  \mathsf{ct}(\ell) = 1 \text{ pour un littéral } \ell,
\]
\[
  \mathsf{ct}\Bigl(\bigwedge_{i=1}^{k} \varphi_i\Bigr) = \prod_{i=1}^{k} \mathsf{ct}(\varphi_i),
  \qquad
  \mathsf{ct}\Bigl(\bigvee_{i=1}^{k} \varphi_i\Bigr) = \sum_{i=1}^{k} \mathsf{ct}(\varphi_i).
\]
La règle multiplicative aux noeuds AND est correcte par
décomposabilité : les enfants ont des variables disjointes, leurs
modèles se combinent indépendamment. La règle additive aux noeuds OR
est correcte par déterminisme : les enfants ont des ensembles de
modèles disjoints, leur somme ne double aucun compte. La mémoisation
sur les identifiants des noeuds donne un parcours en temps $O(|D|)$, où
$|D|$ désigne la taille du DAG (nombre de noeuds plus arêtes).

Appliqué à la racine $\varphi_r(\emptyset, \emptyset)$ du DAG produit
par la compilation BCMS (section \ref{sec:theo-bcms}), $\mathsf{ct}$
livre directement \#SAT$(F)$ : par le Lemme \ref{lem:bcms-root}, les
modèles de $\varphi_r(\emptyset, \emptyset)$ sur $\mathsf{var}(F)$ sont
précisément les affectations satisfaisantes de $F$.

\subsection{CO, VA et ME}
\label{ssec:theo-co-va-me1}

Trois requêtes se déduisent directement du comptage de
sous-section \ref{ssec:theo-ct}.

La consistance (CO) revient à
$\mathsf{ct}(\varphi_r(\emptyset, \emptyset)) > 0$ : le DAG est
consistant si et seulement s'il a au moins un modèle. La décomposabilité
seule suffit pour cette requête, car on s'intéresse à la positivité du
compte et non à sa valeur.

La validité (VA) compare le compte au cardinal de l'espace
d'affectations : $\varphi_r$ est valide si et seulement si
$\mathsf{ct}(\varphi_r(\emptyset, \emptyset)) = 2^n$, où
$n = |\mathsf{var}(F)|$. Cette requête exige déterminisme et lissité
pour que le compte soit exact sur l'espace total des affectations de
$\mathsf{var}(F)$.

L'énumération des modèles (ME) procède top-down par descente dans le
DAG, guidée par le compte $\mathsf{ct}$ déjà calculé en
sous-section \ref{ssec:theo-ct}. À chaque noeud AND, on collecte les
choix de tous les enfants (variables disjointes par décomposabilité).
À chaque noeud OR, on choisit un enfant à compte non nul. Aux feuilles
littérales, on enregistre la valeur du littéral dans l'affectation en
construction. La descente produit un modèle complet en $O(|D|)$.

Le même squelette donne l'énumération de \emph{tous} les modèles : à
chaque noeud OR, au lieu de choisir un enfant on itère sur tous les
enfants à compte non nul, et à chaque noeud AND, le produit cartésien
des sous-énumérations remplace la simple union des choix. Le délai entre
deux modèles consécutifs reste $O(|D|)$, mais le nombre total de
modèles peut être exponentiel. Sur un DAG non lisse, l'énumération
peut produire le même modèle plusieurs fois, et l'opération de
lissage de la section \ref{ssec:theo-trans} y remédie.

\subsection{Transformations internes : conditionnement et lissage}
\label{ssec:theo-trans}

Deux transformations internes sur le d-DNNF préparent le DAG pour les
requêtes CE (section \ref{ssec:theo-ce}) et IM
(section \ref{ssec:theo-im}).

Le \emph{conditionnement} de $\varphi$ par un littéral $\ell$, noté
$\varphi \mid \ell$, remplace chaque occurrence de $\ell$ par $\top$
et chaque occurrence de $\neg \ell$ par $\bot$ dans le DAG, puis
propage les constantes par les règles habituelles
($\top \wedge X \equiv X$, $\bot \wedge X \equiv \bot$,
$\top \vee X \equiv \top$, $\bot \vee X \equiv X$). Pour un terme
$\gamma = \ell_1 \wedge \cdots \wedge \ell_k$, on applique le
conditionnement littéral par littéral. Le résultat $\varphi \mid \gamma$
a pour modèles les affectations
$\tau : \mathsf{var}(\varphi) \setminus \mathsf{var}(\gamma) \to \{0, 1\}$
telles que $\tau \cup \{\ell_i \mapsto 1\}_{i=1}^{k}$ satisfait
$\varphi$.

Le conditionnement préserve la décomposabilité et le déterminisme. Il
ne préserve pas en général la lissité, car la propagation des
constantes peut faire disparaître asymétriquement certaines variables
des sous-DAGs d'un même OR, créant un OR à scopes hétérogènes.

Le \emph{lissage} d'un d-DNNF rétablit la lissité quand elle a été
perdue. Pour chaque noeud OR dont les enfants ont des scopes différents,
on remplace chaque enfant $\psi_i$ par
$\psi_i \wedge \bigwedge_{x \in \mathit{missing}_i} (x \vee \neg x)$,
où $\mathit{missing}_i$ est l'ensemble des variables présentes dans au
moins un autre enfant mais absentes de $\psi_i$. Les termes
$(x \vee \neg x)$ sont identiquement vrais, donc les modèles du DAG
sont inchangés et la lissité est restaurée.

Le lissage est nécessaire chaque fois qu'une requête compare le compte
à une valeur absolue sur l'espace total des affectations (par exemple
le test $\mathsf{ct} = 2^{n-k}$ de IM, section \ref{ssec:theo-im}).
Pour CE en revanche (section \ref{ssec:theo-ce}), le test se réduit
à $\mathsf{ct} = 0$, qui est invariant par lissage : on peut omettre
cette étape.

\subsection{CE : clausal entailment}
\label{ssec:theo-ce}

L'implication clausale teste si une clause
$C = \ell_1 \vee \cdots \vee \ell_k$ est conséquence logique de la
formule $F$ représentée par $\varphi_r(\emptyset, \emptyset)$. La
réduction standard passe par la négation : $F \models C$ équivaut à
dire qu'aucun modèle de $F$ ne satisfait
$\neg C = \neg \ell_1 \wedge \cdots \wedge \neg \ell_k$. Cette
condition s'écrit
\[
  F \models C
  \quad \Longleftrightarrow \quad
  \mathsf{ct}\bigl(\varphi_r(\emptyset, \emptyset) \mid \neg C\bigr) = 0.
\]
On conditionne donc le DAG par le terme $\neg C$
(section \ref{ssec:theo-trans}) puis on calcule le compte du résultat
(section \ref{ssec:theo-ct}). La clause $C$ est conséquence de $F$ si
et seulement si ce compte est nul.

Le test étant fondé sur la nullité du compte, le lissage n'est pas
requis : pour tout d-DNNF $\psi$, $\mathsf{ct}(\psi) = 0$ équivaut à
$\psi \equiv \bot$, indépendamment de la lissité. L'ensemble de la
procédure (conditionnement par $k$ littéraux et comptage) est en
temps polynomial en $|D|$ et $k$.

\subsection{IM : implicant check}
\label{ssec:theo-im}

Le test d'implicant prend en entrée un terme
$\gamma = \ell_1 \wedge \cdots \wedge \ell_k$ portant sur $k$ variables
distinctes, et détermine si toute affectation qui satisfait $\gamma$
satisfait aussi $F$. La réduction passe par le conditionnement et un
test d'égalité sur le compte :
\[
  \gamma \text{ implicant de } F
  \quad \Longleftrightarrow \quad
  \mathsf{ct}(\varphi') = 2^{n-k},
\]
où $n = |\mathsf{var}(F)|$ et $\varphi'$ désigne le lissage de
$\varphi_r(\emptyset, \emptyset) \mid \gamma$. Après conditionnement,
le DAG résiduel encode $F \mid \gamma$ sur les $n - k$ variables
restantes. Le terme $\gamma$ est implicant si et seulement si toutes
les $2^{n - k}$ affectations possibles de ces variables satisfont la
formule, ce qui équivaut à compter exactement $2^{n - k}$ modèles.

À la différence de CE, le test compare le compte à une valeur absolue
sur l'espace total, et non à zéro. Le compte doit donc être exact
relativement à $\mathsf{var}(F) \setminus \mathsf{var}(\gamma)$, ce qui
impose le lissage du DAG résiduel après conditionnement
(section \ref{ssec:theo-trans}). L'enchaînement complet
(conditionnement, lissage, comptage) est en temps polynomial en $|D|$
et $k$.

\chapter{Implémentation}
\label{chap:implementation}


L'implémentation suit le pipeline théorique présenté au
chapitre \ref{chap:theorie} en quatre phases successives : construction
d'une branch decomposition $(T, \delta)$ par l'heuristique GreedyOrder,
parcours bottom-up qui calcule les $\mathsf{PS}'(F_v)$ (Procédure 1),
parcours top-down qui calcule les $\mathsf{PS}'(F_{\overline{v}})$
(Procédure 2), puis programmation dynamique de la Procédure 3 qui
produit conjointement la valeur de MaxSAT, le compte \#SAT et un
DAG d-DNNF structuré au sens de \cite{bova2015}. Ce DAG sert ensuite de
support aux six requêtes polytimes décrites en
section \ref{sec:theo-queries}.

Le code est écrit en C, sans dépendance hors de la bibliothèque
standard. Un langage compilé donne une exécution plus rapide qu'un
langage interprété, choix justifié par les nombreuses itérations de
la programmation dynamique.

% -----------------------------------------------------------------------------
\section{Representation des PS-sets par bitsets 64-bit}
\label{sec:impl-bitsets}
% -----------------------------------------------------------------------------

Saether et al. \cite{saether2015} recommandent de représenter les
PS-sets (section \ref{ssec:theo-cuts-psset}) comme des vecteurs
booléens de longueur $|\mathsf{cla}(F)|$, ce qui rend les opérations
ensemblistes (test d'appartenance, union, intersection) en temps
proportionnel à $m$. Notre implémentation
prolonge ce choix par un packing 64 bits : chaque bitset est un
tableau de \texttt{uint64\_t} de taille $\lceil m/64 \rceil$, et toute
opération ensembliste se réduit à un parcours mot par mot.

La formule $F$ utilise la même représentation : pour chaque variable
$v$, deux bitsets $\mathsf{mask\_pos}[v]$ et $\mathsf{mask\_neg}[v]$
encodent les clauses où $v$ apparaît respectivement en littéral
positif ou négatif.

% -----------------------------------------------------------------------------
\section{Le Binary Trie : deduplication globale des PS-sets}
\label{sec:impl-trie}
% -----------------------------------------------------------------------------

Saether et al. \cite{saether2015} préconisent une structure de trie
binaire pour gérer les PS-sets : un PS-set vu comme une suite de $m$
bits définit un chemin de la racine à une feuille du trie, où chaque
bit choisit l'enfant gauche (clause absente) ou l'enfant droit (clause
présente). À la feuille correspondante, un identifiant entier unique
est attribué la première fois que le PS-set est inséré, et retrouvé
tel quel à toute insertion ultérieure. Insertion et recherche se font
en $O(m)$, et le test d'égalité de deux PS-sets se ramène ensuite à
une comparaison d'entiers en $O(1)$.

La structure concrète est un pool contigu de noeuds, alloué en bloc
et redoublé chaque fois qu'il se remplit. La capacité initiale du
pool vaut $2 m^2$ dès que $m^2 > 10^5$, et $10^5$ sinon ; cette
valeur empirique couvre les instances de notre suite de tests sans
réallocation. Chaque noeud
du pool tient sur trois champs entiers : les deux index d'enfants et
l'identifiant de PS-set (ou $-1$ si le noeud n'est pas terminal).

La déduplication des PS-sets calculés à un noeud $v$ de la branch
decomposition demande, à chaque insertion dans le PS-set local de $v$,
de vérifier si l'identifiant a déjà été ajouté à $v$. Le coût naïf
est $O(|\mathsf{PS}'(F_v)|)$ par insertion. L'implémentation utilise
un tableau \texttt{seen\_array} indexé par identifiant de PS-set, qui
stocke l'identifiant du noeud d'arbre en cours de traitement : le
test \texttt{seen\_array[ps\_id] == node\_id} suffit alors pour
décider en $O(1)$ si le PS-set est déjà présent localement.

% -----------------------------------------------------------------------------
\section{Procedures 1 et 2 : produits cartesiens et masques de coupe}
\label{sec:impl-proc12}
% -----------------------------------------------------------------------------

Les Procédures 1 et 2 transcrivent directement les récurrences des
sections \ref{ssec:theo-proc1} et \ref{ssec:theo-proc2}. Le seul
ressort propre à l'implémentation tient au masque de coupe et à sa
gestion mémoire dans la Procédure 1.

Le masque de coupe est précalculé à la construction de l'arbre. Pour
chaque noeud $v$, $\mathsf{delta\_mask}[v]$ est un bitset sur les
clauses qui matérialise $\mathsf{cla}(F) \cap \delta(v)$, où
$\delta(v)$ est la coupe associée au noeud
(section \ref{ssec:theo-branch-decomp}). Il se calcule bottom-up par
union des masques des enfants. La Procédure 2 utilise directement ce
masque, puisque $\mathsf{cla}(F_{\overline{v}}) = \mathsf{cla}(F) \cap
\delta(v)$. La Procédure 1 calcule à la volée
$\mathsf{mask\_Fv} = \overline{\mathsf{delta\_mask}[v]} \cap
\mathsf{all\_clauses}$ à chaque appel récursif et le libère avant le
retour, pour ne pas accumuler les bitsets temporaires sur la
profondeur de la récursion.

Les cas de base aux feuilles et à la racine sont codés en dur via les
masques $\mathsf{mask\_pos}$ et $\mathsf{mask\_neg}$.

% -----------------------------------------------------------------------------
\section{Procedure 3 : la boucle critique}
\label{sec:impl-proc3}
% -----------------------------------------------------------------------------

La Procédure 3 mène dans une même triple boucle l'agrégat MaxSAT
additif, le décompte \#SAT multiplicatif et la construction du DAG
d-DNNF (section \ref{sec:impl-dag}). L'astuce de Bui-Xuan, Telle et
Vatshelle (section \ref{ssec:theo-proc3}) ramène les six indices
d'une mise à jour de table à trois indices itérés explicitement :
$(C_{c_1}, C_{c_2}, C'_v) \in \mathsf{PS}'(F_{c_1}) \times
\mathsf{PS}'(F_{c_2}) \times \mathsf{PS}'(F_{\overline{v}})$. Les
trois indices restants $C_v$, $C'_{c_1}$, $C'_{c_2}$ sont calculés à
la volée par les identités déjà énoncées en
section \ref{ssec:theo-proc3}, puis insérés dans le trie pour
récupérer leurs identifiants globaux.

Trois bitsets temporaires $\mathsf{temp}_1$, $\mathsf{temp}_2$,
$\mathsf{temp}_3$ sont alloués une fois à l'entrée de la procédure et
réutilisés à chaque triplet pour porter les calculs de $C'_{c_1}$,
$C'_{c_2}$ et $C_v$. Cette pré-allocation évite $\Theta(k^3)$ appels
à \texttt{malloc} par noeud interne.

La boucle itère sur les tailles effectives des PS-sets, sans branche
conditionnelle pour le cas linéaire. Sur un caterpillar
(section \ref{ssec:theo-th23}), l'un des deux PS-sets enfants est de
cardinal au plus deux, donc la triple boucle à un noeud interne
parcourt $O(k^2)$ triplets au lieu de $O(k^3)$, et la borne du
Théorème \ref{thm:dp-linear} est atteinte sans spécialisation du
code.

% -----------------------------------------------------------------------------
\section{Construction du DAG d-DNNF en parallele du DP}
\label{sec:impl-dag}
% -----------------------------------------------------------------------------

La Procédure 3 construit, dans la même boucle que le DP, le DAG
d-DNNF de l'équation \eqref{eq:bcms-construction}. À chaque triplet
$(C_{c_1}, C_{c_2}, C'_v)$ qui met à jour $\mathsf{Tab}_v(C_v, C'_v)$,
un noeud AND est émis et raccroché au noeud OR de la cellule
$(C_v, C'_v)$. Mener le DP et la compilation en même temps épargne
ainsi une seconde passe sur les $\Theta(k^3)$ triplets par noeud
interne. À la racine, $\varphi_r(\emptyset, \emptyset)$ est le DAG
d-DNNF structuré calculant $F$ (Lemme \ref{lem:bcms-correctness}).

Le noeud OR d'une cellule est créé paresseusement au premier AND
ajouté. Une cellule sans aucun AND valide reste vide et n'alloue
aucun noeud, ce qui correspond aux cellules de compte \#SAT nul.
Tous les noeuds du DAG résident dans un pool unique alloué à l'entrée
de la procédure et libéré en bloc à la fin, sans libération
individuelle.

Le pool partage deux singletons pour les constantes $\top$ et $\bot$,
ce qui ramène leur test d'égalité à une comparaison de pointeurs. La
création d'un noeud AND applique alors les simplifications immédiates
$\bot \wedge X \equiv \bot$ et $\top \wedge X \equiv X$ avant toute
allocation, conformément à la propagation des constantes
(section \ref{ssec:theo-bcms-construction}). La taille du DAG produit
reste majorée par la borne $7 k^3 (n+m)$ du Lemme \ref{lem:bcms-size},
et sera mesurée empiriquement au chapitre \ref{chap:resultats}.

% -----------------------------------------------------------------------------
\section{Heuristique GreedyOrder et modes alternatifs}
\label{sec:impl-greedy}
% -----------------------------------------------------------------------------

L'implémentation de GreedyOrder (section \ref{sec:theo-greedy}) tient
sur deux choix d'optimisation. La quantité qui guide la sélection,
$\mathit{Ldegree}(v)$, compte les voisins de $v$ dans $I(F)$ déjà
placés dans l'ordre. À chaque itération, sa mise à jour ne porte que
sur les voisins non choisis du dernier sommet sélectionné, ce qui
évite un recalcul global. La sélection du meilleur candidat et cette
mise à jour incrémentale coûtent chacune $O((n+m)^2)$ au total. À
$\mathit{Ldegree}$ égal entre plusieurs candidats, le sommet de plus
petit degré dans $I(F)$ est retenu (le moins de clauses pour une
variable, le moins de variables pour une clause), conformément au
pseudocode de \cite{saether2015}. Le caterpillar est assemblé au fil
de la sélection : à chaque itération, la feuille du sommet retenu est
greffée à la racine courante de l'arbre en construction.

Deux modes alternatifs sont disponibles via la ligne de commande et
servent de baselines au chapitre \ref{chap:resultats}. Le mode
\texttt{random} fusionne les feuilles deux à deux dans un ordre
aléatoire et produit en général une ps-width très élevée. Le mode
\texttt{linear} construit un caterpillar déterministe selon l'ordre
des variables DIMACS, chaque clause étant greffée juste après sa
variable de plus grand indice. Ce dernier sert de référence pour
mesurer l'apport propre de GreedyOrder par rapport à un ordre brut.

% -----------------------------------------------------------------------------
\section{Requetes polytimes sur d-DNNF : choix d'implementation}
\label{sec:impl-queries}
% -----------------------------------------------------------------------------

Quatre des six requêtes se réduisent à un comptage des modèles ou à
une descente top-down sur le DAG, selon les règles déjà énoncées en
section \ref{ssec:theo-co-va-me1}. La fonction \texttt{count} mémoïse
les comptes par identifiant de noeud, ce qui permet aux quatre
requêtes \texttt{consistency}, \texttt{validity}, \texttt{find\_model}
et \texttt{enumerate} de partager la même table lorsqu'elles sont
enchaînées sur le même DAG.

Les deux dernières requêtes, \texttt{entails} et \texttt{is\_implicant},
s'appuient sur la transformation \texttt{condition}, qui remplace les
occurrences d'un littéral par $\top$ ou $\bot$ selon la valeur souhaitée
et propage les constantes dans le DAG. Le code conserve volontairement
les sous-DAGs $\top$ dans les noeuds OR, au lieu d'appliquer la
simplification $\top \vee X \equiv \top$. Cette non-simplification
préserve les variables apparaissant dans les sous-DAGs $X$, dont le
lissage ultérieur a besoin pour insérer correctement les fragments
$(v \vee \neg v)$ aux endroits manquants \cite{darwiche2002}. La
transformation \texttt{smooth} matérialise ces fragments en exploitant
une table de portées (les variables présentes dans chaque noeud,
section \ref{ssec:theo-properties}), recalculée à chaque appel à
partir de zéro.

\texttt{entails} et \texttt{is\_implicant} composent \texttt{condition}
avec \texttt{count}, et avec \texttt{smooth} pour
\texttt{is\_implicant}, selon les réductions des
sections \ref{ssec:theo-ce} et \ref{ssec:theo-im}.

% -----------------------------------------------------------------------------
\section{Robustesse : gestion memoire et erreurs}
\label{sec:impl-robustness}
% -----------------------------------------------------------------------------

La gestion mémoire suit une discipline d'ownership stricte. Le pool
du DAG (section \ref{sec:impl-dag}) possède tous ses noeuds, le trie
binaire (section \ref{sec:impl-trie}) possède les siens, et le
tableau de PS-sets stocké à chaque noeud $v$ de l'arbre possède les
bitsets qu'il contient. La libération suit un parcours post-ordre de
l'arbre de décomposition : chaque appel libère d'abord les enfants,
puis $\mathsf{PS}'(F_v)$, $\mathsf{PS}'(F_{\overline{v}})$ et le
$\mathsf{delta\_mask}$ du noeud courant. Les tables de DP ne
touchent pas aux noeuds du DAG, qui survivent jusqu'à la libération
en bloc du pool. Les tables elles-mêmes sont stockées comme tableaux
1D indexés en row-major ($i \cdot \mathrm{cols} + j$), ce qui évite
la double indirection d'un tableau 2D classique et favorise la
localité de cache dans les boucles internes de la Procédure 3.

Si une allocation interne échoue dans le trie binaire ou le pool du
DAG, un drapeau est levé et toutes les opérations suivantes
deviennent des no-ops renvoyant des valeurs neutres. L'appelant
détecte ce drapeau et route vers une sortie d'erreur, ce qui évite
un crash silencieux quand la mémoire vient à manquer.


\chapter{Résultats expérimentaux}
\label{chap:resultats}


L'évaluation expérimentale poursuit trois buts. Premièrement, mesurer
le coût réel de la programmation dynamique et de la compilation en
d-DNNF sur des familles d'instances à ps-width contrôlée ou inconnue.
Deuxièmement, valider empiriquement les bornes théoriques du
chapitre \ref{chap:theorie} (taille du DAG, complexité, équivalence
sémantique avec une formule CNF). Troisièmement, situer l'approche
par rapport à un solveur de production (Z3) sur deux tâches
distinctes : la résolution unique de \#SAT/MaxSAT et l'amortissement
de requêtes multiples sur un même DAG compilé.

% -----------------------------------------------------------------------------
\section{Instances testées}
\label{sec:res-instances}
% -----------------------------------------------------------------------------

Le benchmark évalue le pipeline DP+DAG sur six familles d'instances
choisies pour leur structure : trois familles synthétiques tirées du
papier de Saether et al. \cite{saether2015} (type 1, type 2, type 3),
une famille \texttt{random} de $k$-SAT aléatoire qui sert de
référence non structurée, et deux familles fournies par Loïc
Bardi\footnote{Instances \texttt{tseytin} et \texttt{factorization}
fournies par Loïc Bardi, Projet de Bachelor UNIGE 2024.} qui
représentent des circuits booléens réalistes (multiplication
binaire et factorisation entière).

\paragraph{Type 1 : interval ordering libre.}
Le générateur tire $n + m$ intervalles sur la droite réelle. Une
clause $c$ contient une variable $x$ si leurs intervalles se
chevauchent, et la polarité du littéral est tirée uniformément. Le
graphe d'incidence obtenu est un \emph{interval bigraph}, structure
pour laquelle Saether et al. (Corollaire 3) bornent la ps-width par
$\min(m + 1, 2^t)$, avec $t$ la taille maximale des clauses. La
numérotation DIMACS des variables est ensuite permutée
aléatoirement par défaut, ce qui efface la trace de l'interval
ordering : sur la formule permutée, l'ordre DIMACS brut donne une
mauvaise décomposition, et l'heuristique GreedyOrder
(section \ref{sec:theo-greedy}) doit retrouver une décomposition
proche de l'optimale à partir de la structure du graphe. Les
variantes étiquetées \texttt{\_ordered} préservent l'ordre des right
endpoints d'intervalles et servent de baseline contrôle au mode
\texttt{linear} en section \ref{ssec:res-modes}.

\paragraph{Type 2 : interval ordering à taille fixe.}
Type 2 est la variante de type 1 dans laquelle la taille des
clauses est forcée à exactement $t$ (typiquement $t = 3$ ou
$t = 4$), et où chaque intervalle de clause émet quatre clauses sur
les mêmes variables avec des polarités tirées indépendamment. La
borne de Saether reste $\min(m + 1, 2^t)$, mais l'instance est plus
dense et beaucoup plus souvent UNSAT. La même politique de
permutation et de variante \texttt{\_ordered} s'applique.

\paragraph{Type 3 : XOR circulaires.}
Type 3 construit la formule à partir de fonctions XOR à $t$
variables, satisfaites quand le nombre de littéraux vrais est
impair. Les $n$ variables sont disposées en cercle aux positions
$1, \ldots, n$. Étant donné un pas $s \leq t$, le générateur place
les XOR le long du cercle : la $i$-ème couvre les positions
$[i \cdot s, \ldots, i \cdot s + t - 1] \bmod n$, et l'on émet
$n / s$ XOR au total pour faire un tour complet. Quand $s < t$,
deux XOR consécutives partagent $t - s$ variables, ce qui crée une
chaîne de chevauchements le long du cercle. Chaque XOR est ensuite
transcrit en $2^{t-1}$ clauses CNF (celles qui interdisent les
affectations de parité paire des $t$ variables concernées). Le
graphe d'incidence obtenu est la bipartisation d'un graphe d'arcs
circulaires, ce qui borne la ps-width par $\min(m^2 + 1, 4^t)$.

\paragraph{random : $k$-SAT aléatoire.}
La famille \texttt{random} génère des formules $3$-SAT et $4$-SAT
selon le modèle uniforme classique : $m$ clauses tirées
indépendamment, chacune composée de $k$ littéraux distincts à
polarité aléatoire. Les variantes étiquetées \texttt{\_difficile}
fixent le ratio clauses sur variables près du seuil de
satisfiabilité (autour de $4{,}27$ pour $k = 3$), où les instances
sont à la fois denses et incompressibles. Cette famille sert de
référence non structurée : son graphe d'incidence n'a aucune
propriété particulière, et la ps-width attendue est élevée.

\paragraph{tseytin : multiplication binaire $k \times k$.}
Les instances \texttt{tseytin} encodent en CNF la vérification d'une
multiplication binaire $k \times k$, du circuit $3 \times 3$ au
circuit $6 \times 6$ pour les tailles activées dans notre benchmark. Une
variante \texttt{pure} restreint l'encodage aux portes essentielles.
Le graphe d'incidence n'est pas un interval bigraph, et aucune borne
théorique générique ne s'applique : on mesure la ps-width
empiriquement.

\paragraph{factorization : factorisation 5x5.}
Les deux instances \texttt{factorization\_5x5\_464d} et sa variante
simplifiée encodent la factorisation d'un produit binaire fixé sur
cinq bits. Elles correspondent à des instances structurées
non synthétiques (par opposition aux trois familles tirées du
papier), et servent de charge réaliste pour le pipeline.

\begin{table}[htbp]
\centering
\caption{Couverture du benchmark par famille en mode \texttt{greedy}.
La colonne \og\,Échecs\,\fg{} agrège dépassements de la limite
mémoire de 50 GiB par process, crashes et timeouts. Plages $n$, $m$
et ps-width relevées sur les instances réussies. Le mode
\texttt{linear} traite quelques instances supplémentaires
(section \ref{ssec:res-modes}).}
\label{tab:res-instances}
\begin{tabular}{lrrlll}
\toprule
Famille & \# OK & \# Échecs & $n$ & $m$ & ps-width \\
\midrule
type1         & 11 & 3  & 20--200    & 24--213    & 12--21\,000  \\
type2         & 13 & 5  & 25--200    & 88--800    & 14--456      \\
type3         & 19 & 1  & 30--5\,000 & 60--10\,000 & 6--64       \\
random        & 11 & 11 & 7--30      & 20--85     & 32--16\,400  \\
tseytin       & 4  & 2  & 36--72     & 99--220    & 144--8\,192  \\
factorization & 0  & 2  & {}         & {}         & {}           \\
\midrule
Total         & 58 & 24 & & & \\
\bottomrule
\end{tabular}
\end{table}

Le benchmark teste au total 179 paires instance/mode d'arbre en
passe A, dont 110 réussissent : 58 en mode \texttt{greedy}, 46 en
mode \texttt{linear} et 6 en mode \texttt{random}. La
table \ref{tab:res-instances} montre que GreedyOrder couvre presque
entièrement type 3 (19 sur 20), ramène
la ps-width entre $6$ et $64$ y compris pour $n = 5\,000$, et tient
sur l'essentiel de type 1 et type 2. Les échecs se concentrent sur
\texttt{random} (sept dépassements mémoire et quatre crashes
parmi 22 instances tentées) et sur \texttt{factorization}, où
\texttt{greedy} échoue sur les deux instances tandis que
\texttt{linear} les traite avec une ps-width de $24\,576$. Ces
contrastes entre familles structurent les sections suivantes :
\texttt{tseytin} et \texttt{random} concentrent les ps-width
élevées et les pires temps de DP, alors que les trois familles
synthétiques restent dans le régime polynomial visé par le
papier.

Huit instances de très grande taille sont désactivées dans notre
benchmark : quatre type 3 ($n \in \{50\,000, 100\,000\}$ ou $t = 5$
avec $n \in \{5\,000, 10\,000\}$) et quatre tseytin
($k \in \{7, 10\}$). Les activer aurait sensiblement allongé la
durée totale du benchmark sans révéler de comportement
qualitativement nouveau, leur ps-width attendue et leur temps de
DP prolongeant les tendances déjà tracées par les tailles plus
modestes de la même famille.

% -----------------------------------------------------------------------------
\section{Méthodologie expérimentale}
\label{sec:res-methodo}
% -----------------------------------------------------------------------------

Trois axes structurent la méthodologie : la machine d'exécution
(section \ref{ssec:res-machine}), le protocole de mesure en trois
passes complémentaires (section \ref{ssec:res-protocole}) et la
validation des résultats par invariants théoriques
(section \ref{ssec:res-invariants}). Les cas qui s'écartent du
parcours nominal (échecs en exécution, instances désactivées) sont
recensés en section \ref{ssec:res-degeneres}.

\subsection{Machine et environnement}
\label{ssec:res-machine}

Le benchmark s'exécute sur \emph{racer}, une machine du laboratoire de \textbf{Arnaud Casteigts}
équipée de deux processeurs AMD EPYC 9654 (192 cores physiques au
total, 384 threads via SMT2, boost à environ $3{,}7$ GHz) et de
$755$ GiB de mémoire vive, sous Ubuntu $24{,}04$. Le solveur est
compilé avec GCC $13{,}3$ en \texttt{-O3} et la version $4{,}16$ de
Z3 sert de référence externe.

Deux choix de configuration influencent les mesures. Chaque
processus du solveur est borné à $50$ GiB de mémoire résidente, ce
qui transforme les explosions de ps-width en échecs propres plutôt
qu'en saturation de la machine. Pour le chronométrage le plus
isolé (section \ref{ssec:res-protocole}), les workers sont
épinglés sur quatre cores n'appartenant pas au même groupe de
partage de cache L3, ce qui réduit la contention sur les caches
partagés. Les chiffres absolus présentés ensuite doivent être lus
dans ce contexte, alors que les rapports (DP par rapport à Z3,
\texttt{greedy} par rapport à \texttt{linear}) restent plus
robustes au changement de machine.

\subsection{Protocole de mesure en trois passes}
\label{ssec:res-protocole}

Chaque instance traverse trois passes successives, chacune isolant
un type de mesure différent.

La \emph{passe A} est un large parcours parallèle : 64 jobs
simultanés exécutent le solveur sur toutes les instances et tous
les modes d'arbre, en parallèle d'un appel Z3 pour MaxSAT et d'un
appel Z3 pour la satisfiabilité. Elle produit les mesures
structurelles (ps-width, taille du DAG, \#SAT, MaxSAT, temps par
phase) sans isolation CPU stricte. C'est la passe qui permet de
couvrir le plus grand nombre d'instances dans le budget temps imparti.

La \emph{passe B} reprend chaque instance et la rejoue trois fois
sur quatre workers épinglés (section \ref{ssec:res-machine}). Les
quatre temps obtenus sont agrégés en médiane, minimum, maximum et
écart-type. C'est cette passe qui fournit les temps d'exécution
chronométrés rapportés en section \ref{sec:res-dp}.

La \emph{passe C} est plus spécifique : elle mesure le coût des
requêtes décrites en section \ref{sec:theo-queries} sur un DAG
déjà compilé, sans inclure le temps de compilation. Chaque requête
(CO, VA, CT, ME, CE, IM) est
exécutée cinq fois en mémoire dans un même processus, la première
exécution étant écartée pour absorber l'effet de chauffe du cache.
Cette passe alimente les résultats de la section \ref{sec:res-queries}
et l'analyse du break-even N.

Le timeout par instance vaut $1\,800$ s en passe A pour type 1, type 2
et random, et $3\,600$ s pour type 3, tseytin et factorization, qui
peuvent demander plus de temps pour les plus grandes tailles. Les
appels Z3 sont coupés à $600$ s. La durée totale du benchmark est
de $1$ h $50$ min.

\subsection{Validation par invariants théoriques}
\label{ssec:res-invariants}

Huit invariants sont vérifiés après chaque exécution. Ils croisent
les sorties du DP, du DAG compilé et de Z3 pour détecter toute
divergence sémantique ou tout dépassement de borne théorique. La
table \ref{tab:res-invariants} les résume.

\begin{table}[htbp]
\centering
\caption{Invariants vérifiés sur les 164 exécutions réussies (toutes
modes et passes confondues). \og\,WARN\,\fg{} signale un dépassement
\emph{soft} d'une borne théorique générique (sans incidence sur la
correction de l'algorithme). \og\,Skip\,\fg{} compte les instances où
l'invariant ne s'applique pas (par exemple Z3 en timeout pour les
comparaisons croisées).}
\label{tab:res-invariants}
\begin{tabular}{lrrr}
\toprule
Sémantique de l'invariant & \# OK & \# WARN & \# Skip \\
\midrule
MaxSAT(DP) = MaxSAT(Z3)                              & 164 & 0  & 69  \\
\#SAT(DP) = comptage sur le DAG                      & 164 & 0  & 69  \\
SAT(DP) cohérent avec SAT(Z3)                        & 164 & 0  & 69  \\
$|D| \leq 7 k^3 (n + m)$ (Lemme \ref{lem:bcms-size}) & 161 & 0  & 72  \\
ps-width $\leq$ borne théorique de la famille        & 51  & 33 & 149 \\
temps par phase $\geq 0$                             & 164 & 0  & 69  \\
CE(DAG) cohérente avec Z3                            & 54  & 0  & 0   \\
nombre de modèles énumérés = \#SAT                   & 15  & 0  & 39  \\
\bottomrule
\end{tabular}
\end{table}

Aucun invariant ne produit de FAIL hard sur l'ensemble du benchmark.
Les 33 WARN portent uniquement sur l'invariant de ps-width par
rapport à la borne théorique de la famille, et se répartissent en
sept instances type 1, vingt instances type 2 et six instances type
3 où GreedyOrder rate la borne $\min(m+1, 2^t)$ ou
$\min(m^2+1, 4^t)$ du papier \cite{saether2015} sans pour autant
invalider l'algorithme. L'analyse de ces WARN figure
section \ref{ssec:res-pswidth}.

\subsection{Cas dégénérés}
\label{ssec:res-degeneres}

Sur les $82$ paires (instance, mode \texttt{greedy}), $58$
aboutissent et $24$ échouent : $14$ dépassements de la limite
mémoire de $50$ GiB, $8$ segfaults sur des instances très grandes et
$2$ timeouts. Les dépassements mémoire correspondent à des
ps-width effectives au-delà de la dizaine de milliers (familles
\texttt{random\_difficile}, certaines type 1 et type 2), où la
construction des PS-sets explose. Les segfaults se concentrent sur
les deux tailles extrêmes de type 3 ($n = 10\,000$) et sur les
deux instances de factorization en mode \texttt{greedy}. Les
timeouts sont la limite naturelle des plus grands circuits tseytin.

Les sections de résultats portent sur les 58 instances réussies en
mode \texttt{greedy}. Les modes d'échec et leur lien avec la
ps-width sont repris au chapitre \ref{chap:discussion}.
\newpage
% -----------------------------------------------------------------------------
\section{Résultats : programmation dynamique et compilation}
\label{sec:res-dp}
% -----------------------------------------------------------------------------

Cette section regroupe les mesures structurelles et chronométriques
issues des passes A et B (section \ref{ssec:res-protocole}) sur les
$58$ instances réussies en mode \texttt{greedy} : ps-width effective
par famille, temps de résolution comparés à Z3, impact du choix
d'arbre, taille du DAG produit, et corrélation entre la ps-width
mesurée et le nombre de conflits Z3.


\subsection{ps-width observée et bornes théoriques}
\label{ssec:res-pswidth}

La ps-width est le paramètre central des bornes de complexité du
chapitre \ref{chap:theorie} (Théorème \ref{thm:dp-total} :
$O(k^3 m (m+n))$ pour le DP, Lemme \ref{lem:bcms-size} :
$O(k^3 (n+m))$ pour le DAG). La table \ref{tab:res-pswidth}
récapitule les valeurs effectivement produites par GreedyOrder
(section \ref{sec:theo-greedy}) sur les six familles.

\begin{table}[htbp]
\centering
\caption{ps-width observée en mode \texttt{greedy} sur les
instances réussies. La colonne \og\,Borne théorique\,\fg{} reprend
les bornes établies par Saether et al. \cite{saether2015}
(Corollaire 3) pour les familles concernées.}
\label{tab:res-pswidth}
\begin{tabular}{lrrl}
\toprule
Famille & Médiane & Maximum & Borne théorique \\
\midrule
type1   & 32     & 21\,000 & $\min(m + 1,\ 2^t)$ \\
type2   & 184    & 456     & $\min(m + 1,\ 2^t)$ \\
type3   & 8      & 64      & $\min(m^2 + 1,\ 4^t)$ \\
random  & 415    & 16\,400 & (pas de borne) \\
tseytin & 4\,168 & 8\,192  & (pas de borne) \\
\bottomrule
\end{tabular}
\end{table}

Type 3 reste très en dessous de sa borne : les instances $t = 3$
plafonnent à une ps-width de $8$ (borne $4^3 = 64$), et le maximum
observé de $64$ provient des instances $t = 5$ (borne $4^5 = 1\,024$).
GreedyOrder y maintient la ps-width entre $6$ et $64$ y compris pour
$n = 5\,000$. Type 1 et type 2 dépassent en revanche la borne sur
une part importante des instances : c'est l'origine de la majorité
des $33$ alertes \texttt{psw\_within\_family\_bound} signalées en
section \ref{ssec:res-invariants} ($7$ pour type 1 et $20$ pour
type 2). Pour
type 2, la borne nominale $2^t = 8$ avec $t = 3$ est presque
toujours ratée, la médiane atteignant $184$. L'écart vient
principalement de l'heuristique gloutonne, qui ne reconnaît pas
toujours l'interval ordering caché par la permutation aléatoire des
identifiants (section \ref{sec:res-instances}). Les familles
\texttt{random} et \texttt{tseytin}, sans borne théorique
prédéfinie, montent à plusieurs milliers, ce qui annonce les temps
de DP correspondants (section \ref{ssec:res-temps}).

\subsection{Temps DP par famille et comparaison à Z3}
\label{ssec:res-temps}

La table \ref{tab:res-dp-vs-z3} donne les médianes de la passe B
pour le DP en mode \texttt{greedy} et pour Z3 \texttt{Optimize}
sur MaxSAT. La figure \ref{fig:res-dp-vs-z3} reporte la même
mesure par instance.

\begin{table}[htbp]
\centering
\caption{Temps de résolution médian par famille en mode
\texttt{greedy} (passe B) et facteur DP/Z3. Les colonnes
\og\,DP$>$Z3\,\fg{} et \og\,DP$<$Z3\,\fg{} comptent les instances
où DP est respectivement plus lent ou plus rapide que Z3.}
\label{tab:res-dp-vs-z3}
\begin{tabular}{lrrrrr}
\toprule
Famille & DP (ms) & Z3 (ms) & DP/Z3 & DP$>$Z3 & DP$<$Z3 \\
\midrule
type1   & 11{,}9    & 1{,}4 & 8{,}3$\times$       & 10 & 1 \\
type2   & 125{,}7   & 3{,}3 & 38$\times$          & 13 & 0 \\
type3   & 143{,}1   & 4{,}5 & 32$\times$          & 17 & 2 \\
random  & 42{,}3    & 1{,}0 & 41$\times$          & 11 & 0 \\
tseytin & 76\,526{,}1 & 1{,}6 & 47\,734$\times$   & 4  & 0 \\
\bottomrule
\end{tabular}
\end{table}

Z3 résout plus vite que le DP sur $55$ des $58$ instances
comparées. Les trois cas où le DP l'emporte sont
\texttt{type1\_v20\_c25} ($0{,}6$ ms contre $3{,}6$ ms pour Z3) et
les deux variantes permutée et ordonnée de
\texttt{type3\_n30\_t3\_s2} (gain inférieur à $10\,\%$). Tous
sont en deçà de $30$ variables, dans la zone où l'\emph{overhead}
initial de Z3 pèse plus que le coût d'une passe DP. L'écart
médian va d'un facteur $8$ sur type 1 à $47\,734$ sur tseytin, où
le DP s'écroule à cause des ps-width au-delà de $4\,000$
(table \ref{tab:res-pswidth}). Le pire cas absolu est
\texttt{type1\_v200\_c250} (ps-width $21\,000$, $456$ secondes
pour le DP contre $5$ ms pour Z3). Cette infériorité en
single-shot était attendue : Z3 est un solveur industriel
optimisé sur des décennies, alors que le DP poursuit une garantie
de complexité paramétrée (Théorème \ref{thm:dp-total}) qui ne
devient avantageuse que sur des classes structurées à ps-width
réellement bornée. Type 3, où la ps-width reste basse, illustre
cette zone favorable : le DP y prend des temps comparables aux
autres familles malgré $n = 5\,000$ variables, là où Z3 traite
l'instance en quelques millisecondes par effort CDCL standard.

\begin{figure}[H]
\centering
\includegraphics[width=0.75\textwidth]{images/dp_vs_z3_maxsat.png}
\caption{Temps DP en mode \texttt{greedy} (abscisse) par rapport à
Z3 \texttt{Optimize} sur MaxSAT (ordonnée), instance par instance.
La diagonale $y = x$ sépare les instances où le DP est plus
rapide (au-dessus) de celles où Z3 l'emporte (en-dessous). Trois
instances se trouvent au-dessus de la diagonale, toutes de très
petite taille ($n \leq 30$, $m \leq 60$).}
\label{fig:res-dp-vs-z3}
\end{figure}

\subsection{Impact du mode de construction d'arbre}
\label{ssec:res-modes}

La figure \ref{fig:res-greedy-linear} compare les ps-width
obtenues par les deux modes \texttt{greedy} et \texttt{linear} sur
les instances où les deux complètent. GreedyOrder produit
systématiquement une décomposition meilleure ou équivalente, avec
un ratio médian $\mathsf{psw}_{\texttt{linear}} /
\mathsf{psw}_{\texttt{greedy}} = 3{,}1$ et un maximum à $6\,144$.

Le cas le plus extrême est \texttt{type3\_n30\_t3\_s2} :
GreedyOrder atteint $\mathsf{psw} = 8$ (la borne $4^t$ théorique),
alors que \texttt{linear} sur la même instance permutée monte à
$49\,152$. L'ordre des variables DIMACS, arbitraire après la
permutation aléatoire imposée par défaut (section \ref{sec:res-instances}),
ne reflète plus la structure de la formule, et le mode
\texttt{linear} le subit directement. Lorsque la permutation est
désactivée (variantes \texttt{\_ordered}), \texttt{linear}
retrouve des ps-width proches de GreedyOrder, ce qui confirme que
l'apport de l'heuristique est de récupérer un ordre adapté à
partir du seul graphe d'incidence (section \ref{sec:impl-greedy}).

\begin{figure}[H]
\centering
\includegraphics[width=0.75\textwidth]{images/greedy_vs_linear_pswidth.png}
\caption{ps-width obtenue par les modes \texttt{greedy} (vert) et
\texttt{linear} (rouge) sur les 25 instances où l'écart est le
plus marqué, parmi les 40 instances où les deux modes complètent.
Les instances sont classées par ratio
$\mathsf{psw}_{\texttt{linear}} / \mathsf{psw}_{\texttt{greedy}}$
(valeur indiquée à droite de chaque ligne).}
\label{fig:res-greedy-linear}
\end{figure}

\subsection{Corrélation entre ps-width DP et conflits Z3}
\label{ssec:res-correlation}

La figure \ref{fig:res-conflicts-pswidth} confronte la ps-width
mesurée par le DP au nombre de conflits internes consommés par
Z3 lors de sa résolution.

Les deux quantités ne sont pas alignées. Les instances tseytin,
dont la ps-width grimpe entre $144$ et $8\,192$, sont résolues par
Z3 avec zéro ou deux conflits dans tous les cas observés : Z3 ne
fait quasiment pas de backtrack, sa résolution est portée par la
propagation unitaire seule. À l'inverse, type 2 montre une
ps-width modérée (médiane $184$) pour une médiane de huit
conflits, le maximum observé restant en deçà de la centaine. Aucun
cas du benchmark ne combine ps-width faible et nombre de
conflits élevé, configuration qui aurait pu jouer en faveur du DP
contre Z3. Ces observations confirment que la ps-width et l'effort
CDCL mesurent des notions de difficulté différentes : la première
est une propriété structurelle du graphe d'incidence
(section \ref{ssec:theo-pswidth-def}), alors que la seconde reflète
le comportement local d'un solveur d'inférence.

\begin{figure}[H]
\centering
\includegraphics[width=0.75\textwidth]{images/z3_conflicts_vs_pswidth.png}
\caption{Conflits Z3 par rapport à la ps-width effective (mode
\texttt{greedy}), instance par instance. Chaque famille est
représentée par une couleur distincte.}
\label{fig:res-conflicts-pswidth}
\end{figure}

% -----------------------------------------------------------------------------
\section{Résultats : requêtes sur DAG compilé}
\label{sec:res-queries}
% -----------------------------------------------------------------------------

Une fois le DAG compilé, les requêtes polytimes de la
section \ref{sec:theo-queries} s'exécutent par simple parcours, sans
revenir au DP. La passe C mesure leur coût propre sur le DAG déjà
construit (section \ref{ssec:res-protocole}), ce qui permet de
quantifier le seuil à partir duquel ce coût de compilation est
amorti par un usage répété.

\subsection{Coût par requête}
\label{ssec:res-cout-requete}

La table \ref{tab:res-queries} donne le temps médian de chaque
requête sur l'ensemble des DAG compilés, ainsi que son coût ramené
à l'arête du DAG.

\begin{table}[H]
\centering
\caption{Coût médian des requêtes sur DAG compilé (passe C, médiane
sur 5 répétitions après warm-up), sur les 39 instances à DAG non
vide. Le coût par arête rapporte le temps à la taille du DAG. La ligne \og\,tous les modèles\,\fg{} est
mesurée sur les 15 instances dont l'ensemble des modèles reste
énumérable (médiane de 47 modèles). Sur les autres, ce nombre
dépasse les centaines de milliers et l'énumération complète n'est
pas chronométrée.}
\label{tab:res-queries}
\begin{tabular}{lrr}
\toprule
Requête & Médiane (µs) & Coût par arête (ns) \\
\midrule
CO (consistance)            & 42{,}8  & 10{,}6  \\
VA (validité)               & 44{,}0  & 10{,}5  \\
CT (comptage \#SAT)         & 42{,}6  & 10{,}5  \\
CE (implication clausale)   & 295{,}6 & 66{,}9  \\
IM (test d'implicant)       & 155{,}4 & 161{,}2 \\
ME, un modèle               & 50{,}0  & 11{,}2  \\
ME, tous les modèles        & 47{,}5  & {}      \\
\bottomrule
\end{tabular}
\end{table}

\newpage
Les requêtes CO, VA et CT coûtent entre $43$ et $44$ µs, soit
environ $10$ ns par arête, ce qui reflète leur structure commune :
un parcours unique du DAG en temps $O(|D|)$
(section \ref{ssec:theo-ct}). CE et IM sont plus chères
(respectivement $296$ et $155$ µs) car elles conditionnent d'abord
le DAG, et IM le lisse ensuite, deux transformations qui le
parcourent avant le comptage (section \ref{sec:impl-queries}).

Les deux variantes de ME se distinguent par le traitement des
noeuds OR. La recherche d'un modèle ne descend qu'une branche par
OR et s'arrête au premier modèle trouvé ($50$ µs). L'énumération
parcourt au contraire toutes les branches pour produire l'ensemble
des modèles. Sur les $15$ instances où cet ensemble reste
énumérable, elle prend $47{,}5$ µs, du même ordre que la recherche
d'un seul modèle, car ces instances ont peu de modèles (médiane
$47$) et l'énumération ne précalcule pas la table des comptes
nécessaire à la descente unique.

\subsection{Comparaison CO vs Z3 SAT sur la même question}
\label{ssec:res-co-z3}

Parmi les six requêtes, seule la consistance (CO) pose à Z3 une
question identique : la formule est-elle satisfaisable. Les autres
(VA, CT, CE, IM, ME) répondent à des questions que Z3 ne traite pas
nativement et ne donnent donc pas lieu à une comparaison directe.
La figure \ref{fig:res-co-z3} confronte uniquement CO sur DAG
compilé à une décision SAT de Z3, sur une instance représentative
de chaque famille.

\begin{figure}[H]
\centering
\includegraphics[width=0.8\textwidth]{images/co_vs_z3_per_instance.png}
\caption{Temps de la requête de consistance (CO) sur DAG compilé
(bleu) face à une décision SAT de Z3 (rouge), sur une instance
représentative par famille. Le facteur d'accélération est indiqué à
droite. Échelle logarithmique en abscisse.}
\label{fig:res-co-z3}
\end{figure}

Sur ces instances, CO sur DAG va de $4$ à $71$ fois plus vite qu'un
appel SAT de Z3. L'écart est le plus net sur \texttt{random\_k3\_v8\_c20}
($\times 71$) et \texttt{type2\_v100\_c400\_t4\_ordered}
($\times 50$), et le plus mince sur \texttt{type3\_n90\_t5\_s3\_ordered}
($\times 4$). Cet avantage suppose que le DAG est déjà construit. Il
faut aussi noter que Z3 est ici relancé à froid pour chaque instance,
sans réutilisation d'état entre requêtes successives, et un solveur
incrémental réduirait l'écart. Cette nuance, et la mise en balance
avec le coût de compilation, sont reprises au
chapitre \ref{chap:discussion}.

\subsection{Break-even N}
\label{ssec:res-breakeven}

La requête isolée est rapide, mais la compilation du DAG ne l'est
pas (section \ref{sec:res-dp}). L'approche n'est rentable que si le
DAG sert à plusieurs requêtes. Le seuil de rentabilité $N^*$ est le
nombre de requêtes au-delà duquel compiler le DAG puis poser $N$
requêtes coûte moins cher que $N$ appels Z3 indépendants. La
figure \ref{fig:res-breakeven} trace les courbes de coût (panneau
gauche) et la distribution de $N^*$ sur l'ensemble des instances
(panneau droit).

\begin{figure}[H]
\centering
\includegraphics[width=\textwidth]{images/breakeven_n.png}
\caption{À gauche : coût total du pipeline DP+requêtes en fonction
du nombre $N$ de requêtes (trait plein), comparé à $N$ appels Z3
froids (pointillés), pour une instance représentative par famille.
À droite : distribution cumulée du seuil $N^*$ sur les instances où
il est fini, médiane à $120$.}
\label{fig:res-breakeven}
\end{figure}

La médiane de $N^*$ vaut $120$, avec un premier quartile à $19$ et
un troisième quartile à $370$. Autrement dit, dès qu'une même
formule fait l'objet d'environ $120$ requêtes, le pipeline DP+DAG
devient en moyenne plus économique que des appels Z3 répétés. Le
seuil est fini sur $36$ des $39$ instances mesurées. Les trois
exceptions, où aucun nombre raisonnable de requêtes n'amortit la
compilation, sont les instances à ps-width élevée dont le DAG est
le plus coûteux à produire. L'amortissement multi-requêtes est donc
le cas d'usage où la compilation en d-DNNF prend tout son sens,
par opposition à la résolution unique où Z3 domine
(section \ref{ssec:res-temps}).


\chapter{Discussion}
\label{chap:discussion}

% -----------------------------------------------------------------------------
\section{Du solveur au compilateur de connaissances}
\label{sec:disc-pivot}
% -----------------------------------------------------------------------------

Pris comme solveur de \#SAT et de MaxSAT à coup unique,
l'algorithme de programmation dynamique de Saether et
al. \cite{saether2015} ne rivalise pas avec un solveur de
production. La section \ref{ssec:res-temps} chiffre l'écart : Z3
résout plus vite sur la quasi-totalité des instances comparées, et
le facteur de ralentissement du DP atteint environ $47\,734$ sur
les circuits tseytin. Cette infériorité était attendue. Z3
capitalise des décennies d'optimisation de la recherche CDCL, là
où la programmation dynamique poursuit une garantie de complexité
paramétrée par la ps-width, qui ne devient favorable que sur des
familles réellement structurées.

L'intérêt de l'approche ne tient donc pas à sa vitesse de
résolution. La même exécution qui calcule \#SAT et MaxSAT construit
un DAG d-DNNF structuré (section \ref{ssec:theo-bcms-construction}),
et ce DAG est un objet réutilisable : une fois compilé, il répond à
toute une famille de requêtes par simple parcours, sans relancer la
programmation dynamique (section \ref{sec:res-queries}). La carte de
compilation de Darwiche et Marquis \cite{darwiche2002} garantit que
le d-DNNF rend polynomiales les six requêtes exploitées dans ce
travail : consistance, validité, comptage des modèles, implication
clausale, test d'implicant et énumération
(section \ref{sec:theo-queries}). La contribution de ce mémoire
n'est ainsi pas un solveur plus rapide, mais un compilateur de
connaissances.

La compilation a toutefois un coût, celui de la programmation
dynamique elle-même (section \ref{sec:res-dp}), alors qu'une requête
isolée sur le DAG est presque gratuite. L'approche n'est rentable
que si le même DAG sert à plusieurs requêtes. La
section \ref{ssec:res-breakeven} quantifie ce seuil : sur la médiane
des instances mesurées, compiler le DAG puis poser $N^* = 120$
requêtes coûte autant que $120$ appels Z3 indépendants, et
l'avantage passe au compilateur au-delà. C'est dans cet usage
répété, et non dans la résolution unique, que la compilation en
d-DNNF prend sa valeur.

% -----------------------------------------------------------------------------
\section{ps-width et effort CDCL : des difficultés orthogonales}
\label{sec:disc-pswidth-cdcl}
% -----------------------------------------------------------------------------

La ps-width et l'effort de recherche d'un solveur CDCL mesurent
deux difficultés sans rapport l'une avec l'autre. La
section \ref{ssec:res-correlation} le met en évidence : les
instances tseytin atteignent une ps-width allant jusqu'à $8\,192$,
là où Z3 les résout avec zéro à deux conflits. Une ps-width élevée
n'entraîne donc pas une résolution coûteuse pour Z3, et rien dans
nos mesures ne suggère la relation inverse. L'écart tient à la
nature des deux quantités. La ps-width est une propriété
structurelle du graphe d'incidence de la formule
(section \ref{ssec:theo-pswidth-def}) et borne la largeur des états
que la programmation dynamique maintient le long de la
décomposition. Le nombre de conflits, lui, mesure l'effort de
recherche que Z3 consent au fil de sa résolution. Aucune des deux
ne se déduit de l'autre.

Cette indépendance pèse directement sur la comparaison de la
section \ref{sec:disc-pivot}. Le seul régime qui aurait avantagé la
programmation dynamique, une ps-width basse couplée à un grand
nombre de conflits, n'apparaît dans aucune instance testée, tandis
que les familles à ps-width élevée que le DP peine à traiter sont
justement celles que Z3 résout sans effort. Aucun des deux
paradigmes ne domine l'autre en général : chacun capture une forme
de difficulté que l'autre ignore.

% -----------------------------------------------------------------------------
\section{GreedyOrder, maillon faible et principal levier}
\label{sec:disc-greedy}
% -----------------------------------------------------------------------------

La qualité des résultats du DP dépend entièrement de la branch
decomposition fournie en entrée, et celle-ci provient de
l'heuristique GreedyOrder, que Saether et al. \cite{saether2015}
décrivent eux-mêmes comme \og\,very simple\,\fg{}. Sur les familles
structurées, cette simplicité se paie. Type 2 illustre le cas le
plus net : sa borne théorique vaut $2^t = 8$ pour $t = 3$, alors
que GreedyOrder y produit une ps-width de médiane $184$
(section \ref{ssec:res-pswidth}). Au total, $33$ instances
franchissent la borne de leur famille sans que l'algorithme en soit
pour autant invalidé (section \ref{ssec:res-invariants}).

Ce dépassement est le point le plus coûteux de tout le pipeline. La
complexité du DP croît comme le cube de la ps-width
(Théorème \ref{thm:dp-total}), si bien que diviser la ps-width par
deux divise le temps de résolution par environ huit. Aucun autre
paramètre n'a un effet aussi marqué sur le coût. Or l'écart à la
borne ne vient pas de la formule, mais de l'ordre dans lequel la
décomposition est construite. La numérotation des variables est
permutée aléatoirement par défaut, ce qui masque l'interval
ordering sous-jacent, et l'heuristique ne le reconstitue pas
toujours à partir du seul graphe d'incidence
(section \ref{ssec:res-modes}).

Deux voies pourraient resserrer cet écart. La première est exacte :
pour les formules dont le graphe d'incidence est un \emph{interval
bigraph}, ce qui couvre les familles type 1 et type 2, le
Théorème 6 de Saether et al. \cite{saether2015} fournit en temps
polynomial une décomposition linéaire atteignant la borne
$\min(m + 1,\ 2^t)$. Reconnaître ces formules et leur appliquer
cette construction dédiée remplacerait l'heuristique par une
décomposition optimale là où elle échoue le plus. La seconde est
plus générale : les heuristiques de calcul de décompositions
arborescentes, éprouvées de longue date pour minimiser la
treewidth, pourraient être transposées à la recherche d'une branch
decomposition de faible ps-width. Dans les deux cas, le constat
reste le même. L'écart aux bornes du papier tient à l'heuristique
de décomposition, non à l'algorithme de programmation dynamique,
dont aucune exécution n'a produit de résultat incorrect
(section \ref{ssec:res-invariants}).

% -----------------------------------------------------------------------------
\section{Limites de l'évaluation et améliorations possibles}
\label{sec:disc-limites}
% -----------------------------------------------------------------------------

Plusieurs choix du protocole limitent la portée de la comparaison
avec Z3, et tous flattent notre approche. Z3 est relancé à froid à
chaque appel, sans réutilisation de son état interne d'une requête à
la suivante. L'avantage mesuré de la consistance sur DAG
(section \ref{ssec:res-co-z3}) en ressort surévalué, puisqu'un usage
incrémental de Z3 amortirait une partie de son coût de démarrage. Une
comparaison à un Z3 incrémental, sur des séquences de requêtes
proches, en donnerait une image plus juste. Plus fondamentalement, Z3
est un solveur SMT généraliste, ni un solveur de \#SAT dédié, ni un
compilateur de connaissances. Les solveurs de comptage spécialisés
comme sharpSAT \cite{thurley2006} et les compilateurs comme
d4 \cite{lagniez2017} ou ganak offriraient une comparaison plus
exigeante, que nous n'avons pas conduite, et les intégrer comme
références situerait l'approche par rapport à l'état de l'art réel du
domaine.

Deux autres réserves portent sur la portée des conclusions plutôt que
sur la comparaison. Le benchmark ne retient que les instances menées à
terme (section \ref{ssec:res-degeneres}). Les plus difficiles
dépassent la limite mémoire imposée et n'apparaissent pas dans les
statistiques, qui décrivent donc le comportement du DP sur les
instances faisables et sous-représentent son pire cas. La borne de
taille $7 k^3 (n + m)$ (Lemme \ref{lem:bcms-size}) appelle une
réserve du même ordre. Elle est respectée sur toutes les instances
vérifiées (section \ref{ssec:res-invariants}), mais la taille réelle
du DAG lui reste inférieure d'environ six ordres de grandeur en
médiane. La borne garde sa valeur de garantie théorique et ne sert à
rien pour dimensionner la mémoire d'une exécution concrète.

Ces réserves dessinent les prolongements naturels du côté de
l'évaluation : intégrer un Z3 incrémental et des outils de \#SAT
spécialisés comme références, et repousser la limite mémoire pour
mesurer le DP là où il échoue aujourd'hui. Avec les pistes
algorithmiques de la section \ref{sec:disc-greedy}, elles tracent le
programme d'une suite à ce travail, centrée sur la qualité de la
compilation plutôt que sur la vitesse de résolution.

\chapter{Conclusion}
\label{chap:conclusion}

% -----------------------------------------------------------------------------
\section{Bilan et réponse à la question de recherche}
\label{sec:concl-bilan}
% -----------------------------------------------------------------------------

Ce mémoire visait deux objectifs complémentaires
(chapitre \ref{chap:intro}), l'un d'implémentation, l'autre
d'extension. Le premier est atteint : l'algorithme de programmation
dynamique de Saether et al. \cite{saether2015} est implémenté de bout
en bout en C, depuis la lecture des formules CNF jusqu'au calcul
conjoint de \#SAT et de MaxSAT, en passant par la construction d'une
branch decomposition et les trois procédures du papier. Le second
l'est aussi : la programmation dynamique est instrumentée à la
manière de Bova, Capelli, Mengel et Slivovsky \cite{bova2015} pour
produire, au cours de la même exécution, un DAG d-DNNF structuré, sur
lequel les six requêtes que la carte de compilation de Darwiche et
Marquis \cite{darwiche2002} rend polynomiales ont toutes été
réalisées. Une évaluation empirique sur les familles synthétiques à
ps-width contrôlée, les formules aléatoires et les circuits de
multiplication accompagne ces deux volets
(chapitre \ref{chap:resultats}).

La cohérence de l'ensemble a été contrôlée par une batterie
d'invariants théoriques (section \ref{ssec:res-invariants}). Sur
toutes les instances vérifiées, la valeur de MaxSAT et le comptage des
modèles coïncident avec ceux de Z3, le compte effectué sur le DAG
redonne le \#SAT de la programmation dynamique, et la taille du DAG
respecte sa borne théorique, sans aucun échec. Théorie
(chapitre \ref{chap:theorie}) et implémentation
(chapitre \ref{chap:implementation}) concordent ainsi sur tout le
benchmark.

L'introduction posait une question : dans quels régimes la
décomposition structurelle apporte-t-elle un avantage ? La réponse est
nuancée. Pour résoudre une formule une seule fois, elle n'en apporte
aucun, et un solveur de production traite plus vite la quasi-totalité
des instances (section \ref{ssec:res-temps}). L'avantage se loge
ailleurs, dans la réutilisation du DAG compilé : dès qu'une même
formule fait l'objet d'environ $N^* = 120$ requêtes, le coût de
compilation est amorti et le pipeline devient le plus économique
(section \ref{ssec:res-breakeven}). La décomposition structurelle ne
livre donc pas un meilleur solveur, mais un compilateur de
connaissances dont la valeur tient à l'usage répété
(section \ref{sec:disc-pivot}).

% -----------------------------------------------------------------------------
\section{Perspectives}
\label{sec:concl-perspectives}
% -----------------------------------------------------------------------------

La portée de ce résultat dépasse le cadre du benchmark. Un DAG
compilé une fois sert de support à un raisonnement répété sur la même
formule, qu'il s'agisse de compter ses modèles, d'en énumérer ou de
tester des implications, sans jamais relancer la programmation
dynamique (section \ref{sec:disc-pivot}). Les usages évoqués en
introduction (chapitre \ref{chap:intro}), de l'inférence probabiliste
exacte à la vérification quantitative, relèvent précisément de ce
schéma, où une même théorie est interrogée un grand nombre de fois.

Le seuil de rentabilité mesuré (section \ref{ssec:res-breakeven})
cerne le contexte où la compilation trouve son intérêt :
l'interrogation répétée d'une formule, et non sa résolution
ponctuelle. C'est là, plus que dans la course à la vitesse face aux
solveurs de production, que la programmation dynamique sur branch
decomposition et la compilation en d-DNNF prennent leur sens.

% =============================================================

\begin{thebibliography}{99}

\bibitem{cook1971}
  S. A. Cook.
  \newblock The complexity of theorem-proving procedures.
  \newblock In \emph{Proc.\ 3rd Annual ACM Symposium on Theory of Computing (STOC)}, pp. 151--158, 1971.

\bibitem{karp1972}
  R. M. Karp.
  \newblock Reducibility among combinatorial problems.
  \newblock In R. E. Miller and J. W. Thatcher (eds.), \emph{Complexity of Computer Computations}, pp. 85--103. Plenum Press, 1972.

\bibitem{valiant1979}
  L. G. Valiant.
  \newblock The complexity of enumeration and reliability problems.
  \newblock \emph{SIAM Journal on Computing}, 8(3):410--421, 1979.

\bibitem{moskewicz2001}
  M. W. Moskewicz, C. F. Madigan, Y. Zhao, L. Zhang, and S. Malik.
  \newblock Chaff: Engineering an efficient SAT solver.
  \newblock In \emph{Proc.\ 38th Design Automation Conference (DAC)}, pp. 530--535, 2001.

\bibitem{een2003}
  N. Eén and N. Sörensson.
  \newblock An extensible SAT-solver.
  \newblock In \emph{Proc.\ 6th International Conference on Theory and Applications of Satisfiability Testing (SAT)}, LNCS 2919, pp. 502--518. Springer, 2003.

\bibitem{marques1999}
  J. P. Marques-Silva and K. A. Sakallah.
  \newblock GRASP: A search algorithm for propositional satisfiability.
  \newblock \emph{IEEE Transactions on Computers}, 48(5):506--521, 1999.

\bibitem{thurley2006}
  M. Thurley.
  \newblock sharpSAT -- Counting models with advanced component caching and implicit BCP.
  \newblock In \emph{Proc.\ 9th International Conference on Theory and Applications of Satisfiability Testing (SAT)}, LNCS 4121, pp. 424--429. Springer, 2006.

\bibitem{lagniez2017}
  J.-M. Lagniez and P. Marquis.
  \newblock An improved decision-DNNF compiler.
  \newblock In \emph{Proc.\ 26th International Joint Conference on Artificial Intelligence (IJCAI)}, pp. 667--673, 2017.

\bibitem{moura2008}
  L. de Moura and N. Bjørner.
  \newblock Z3: An efficient SMT solver.
  \newblock In \emph{Proc.\ 14th International Conference on Tools and Algorithms for the Construction and Analysis of Systems (TACAS)}, LNCS 4963, pp. 337--340. Springer, 2008.

\bibitem{darwiche2002}
  A. Darwiche and P. Marquis.
  \newblock A knowledge compilation map.
  \newblock \emph{Journal of Artificial Intelligence Research (JAIR)}, 17:229--264, 2002.

\bibitem{darwiche2011}
  A. Darwiche.
  \newblock SDD: A new canonical representation of propositional knowledge bases.
  \newblock In \emph{Proc.\ 22nd International Joint Conference on Artificial Intelligence (IJCAI)}, pp. 819--826, 2011.

\bibitem{bryant1986}
  R. E. Bryant.
  \newblock Graph-based algorithms for Boolean function manipulation.
  \newblock \emph{IEEE Transactions on Computers}, C-35(8):677--691, 1986.

\bibitem{samer2010}
  M. Samer and S. Szeider.
  \newblock Algorithms for propositional model counting.
  \newblock \emph{Journal of Discrete Algorithms}, 8(1):50--64, 2010.

\bibitem{slivovsky2020}
  F. Slivovsky and S. Szeider.
  \newblock A faster algorithm for propositional model counting parameterized by incidence treewidth.
  \newblock In \emph{Proc.\ 23rd International Conference on Theory and Applications of Satisfiability Testing (SAT)}, LNCS 12178, pp. 267--276. Springer, 2020.

\bibitem{vatshelle2012}
  M. Vatshelle.
  \newblock \emph{New Width Parameters of Graphs}.
  \newblock PhD thesis, University of Bergen, 2012.

\bibitem{bergougnoux2025}
  B. Bergougnoux, É. Bonnet, and J. Duron.
  \newblock Mim-width is paraNP-complete.
  \newblock In \emph{Proc.\ 52nd International Colloquium on Automata, Languages, and Programming (ICALP 2025)}, LIPIcs, 2025.
  \newblock DOI: 10.4230/LIPIcs.ICALP.2025.25.

\bibitem{bova2015}
  S. Bova, F. Capelli, S. Mengel, and F. Slivovsky.
  \newblock On compiling CNFs into structured deterministic DNNFs.
  \newblock In \emph{Proc.\ 18th International Conference on Theory and Applications of Satisfiability Testing (SAT)}, LNCS 9340, pp. 199--214. Springer, 2015.

\bibitem{saether2015}
  S. H. Saether, J. A. Telle, and M. Vatshelle.
  \newblock Solving \#SAT and MaxSAT by dynamic programming.
  \newblock \emph{Journal of Artificial Intelligence Research (JAIR)}, 54:59--82, 2015.

\end{thebibliography}

% Déclaration d'utilisation de l'IA
\chapter*{Déclaration d'utilisation de l'IA}
\addcontentsline{toc}{chapter}{Déclaration d'utilisation de l'IA}

Je déclare avoir eu recours à des assistants d'intelligence
artificielle selon les modalités suivantes :

\begin{itemize}
  \item \textbf{Recherche et théorie} : l'IA a servi d'aide à la
  recherche bibliographique et à la consolidation de ma compréhension
  des concepts abordés dans ce travail. \\

  \item \textbf{Aide au développement} : l'IA a été utilisée comme
  partenaire de travail pour le débogage, l'optimisation, la
  suggestion de structures modulaires pour l'implémentation et la
  réalisation de figures soignées présentant les résultats.
  L'intégralité du code produit demeure néanmoins mon propre travail. \\

  \item \textbf{Rédaction et mise en forme} : l'IA m'a aidé pour la
  correction orthographique et grammaticale du texte et pour la
  reformulation de certains passages, afin d'en améliorer la clarté.
\end{itemize}

\clearpage

\end{document}